% !TeX program = lualatex
\documentclass[11pt,a4paper]{ltjsarticle}
\usepackage[margin=25mm,headheight=17pt]{geometry}
\usepackage[haranoaji]{luatexja-preset}
\usepackage{amsmath,amssymb,amsthm,mathtools,mathrsfs,booktabs,array,enumitem,longtable}
\usepackage[unicode,hidelinks]{hyperref}
\usepackage{fancyhdr}
\pagestyle{fancy}\fancyhf{}\fancyfoot[C]{\thepage}
\fancyhead[L]{\small\nouppercase{\leftmark}}
\renewcommand{\sectionmark}[1]{\markboth{#1}{}}
\setlength{\emergencystretch}{3em}\allowdisplaybreaks[2]
\setlength{\parskip}{.15em}
\newtheoremstyle{jp}{.6em}{.6em}{\normalfont}{}{\bfseries}{.}{.5em}{}
\theoremstyle{jp}
\newtheorem{theorem}{定理}[section]
\newtheorem{proposition}[theorem]{命題}
\newtheorem{lemma}[theorem]{補題}
\newtheorem{corollary}[theorem]{系}
\newtheorem{definition}[theorem]{定義}
\newtheorem{remark}[theorem]{注意}
\newtheorem{example}[theorem]{例}
\newtheorem{question}[theorem]{問題}
\renewcommand{\proofname}{証明}
\newcommand{\kk}{\mathbb C}\newcommand{\PP}{\mathbb P}
\newcommand{\OO}{\mathcal O}\newcommand{\A}{\mathcal A}
\newcommand{\B}{B_{\mathrm{VK}}}\newcommand{\Blo}{B_{p,q}}
\newcommand{\HH}{\operatorname{HH}}\newcommand{\Hom}{\operatorname{Hom}}
\newcommand{\End}{\operatorname{End}}\newcommand{\Sym}{\operatorname{Sym}}
\newcommand{\Perf}{\operatorname{Perf}}\newcommand{\Der}{\operatorname{Der}}
\newcommand{\Inn}{\operatorname{Inn}}\newcommand{\gr}{\operatorname{gr}}
\newcommand{\rk}{\operatorname{rank}}\newcommand{\Spec}{\operatorname{Spec}}
\newcommand{\id}{\operatorname{id}}\newcommand{\Ext}{\operatorname{Ext}}
\newcommand{\Tor}{\operatorname{Tor}}\newcommand{\Ann}{\operatorname{Ann}}
\newcommand{\coker}{\operatorname{coker}}\newcommand{\Supp}{\operatorname{Supp}}
\newcommand{\pd}{\operatorname{pd}}\newcommand{\depth}{\operatorname{depth}}
\newcommand{\Frac}{\operatorname{Frac}}\newcommand{\Aut}{\operatorname{Aut}}
\newcommand{\Out}{\operatorname{Out}}\newcommand{\GL}{\operatorname{GL}}
\newcommand{\PGL}{\operatorname{PGL}}\newcommand{\Hilb}{\operatorname{Hilb}}
\newcommand{\dashsim}{\mathrel{\overset{\mathrm{bir}}{\sim}}}
\setcounter{tocdepth}{1}

\hypersetup{pdftitle={四次Prymデータに付随する線織面と種数二十K3の境界周期},pdfauthor={}}
\title{四次Prymデータに付随する線織面と\\種数二十K3の境界周期}
\author{}\date{2026年10月29日\quad 第30版}
\begin{document}
\maketitle
\begin{abstract}
自己反対Fano72型関係代数の署名 $(2,4)$ の八次元成分を、四次Prymデータから調べる。
種数三Prymモジュライへの写像は支配的で、一般幾何学的ファイバーは
二次曲線netのHermite曲線 $E$ 上の線織面と双有理である。
逆Gram構成から生じる一次行列式表示の線束を $\theta\in\operatorname{Pic}^0(E)$ とすると、
等方Grassmann曲面の主成分は $\mathbb P_E(\mathcal O\oplus\theta)$、
残余符号商は $\mathbb P_E(\mathcal O\oplus\theta^2)$ となり、商写像はファイバー上の二乗写像である。
閉包はこの主成分と二つの射影平面の単純正規交差であり、
一般摂動と十八普通二重点の射影的小解消は、種数二十K3のType II Kulikovモデルを与える。
両端の有理曲面の反標準束は、$\theta$ が有限位数であるとき、かつそのときに限り半豊富である。
幾何学的一般点では $\theta$ は無限位数である。
さらに、偏極に直交する極限整数格子の判別式は三十八、根系は $A_8\oplus A_8$ であり、
モノドロミーの対数の最大整数因子は二である。
さらに、二つの九点配置を独立に動かす平滑化切断の全射性を証明する。
十八点の配置から偏極付き境界周期への写像は、対角的な $E[3]$ を核とする次数九の同種写像であり、
元の逆Gram族はこの境界型の十八次元の周期空間を支配する。
\end{abstract}
\begingroup\small\setlength{\parskip}{0pt}\tableofcontents\endgroup
\clearpage

\section{序論と主結果}

全自己反対方程式は、関係行列の二つの固有ブロックに対称Gram形式を与え、
異符号ブロックの間に一次syzygyを生む。
署名 $(2,4)$ の場合には、これらの式から四次曲線と非自明二捩れ線束の対、
その二次曲線netのHermite曲線、六次元上の三つの交代形式が自然に現れる。
本稿は、この対応から得る線織面と、等方Grassmann曲面の境界・K3平滑化を扱う。
$\mathcal R_g$ は曲線と非自明二捩れ線束の対の粗モジュライである。
Prym多様体だけを指定するモジュライとは区別する。

\begin{theorem}[四次PrymファイバーとそのK3退化]
\label{thm:main-prym-geometry-v29}
標数零の代数閉体上で、指定反対合を保持した署名 $(2,4)$ の有理商について、次が成立する。
\begin{enumerate}
\item 八次元商からの写像 $\Phi_3:\mathfrak Z_3\dashrightarrow\mathcal R_3$ は支配的である。
一般幾何学的ファイバーはnetのHermite曲線 $E$ 上の射影直線束と双有理であり、
\[
 2\det G\operatorname{Pf}\Omega(z)=H_N(z)
\]
という普遍恒等式がこの曲線を同定する。
\item 二つの共通等方三平面を順序付け、
$\Omega=\left(\begin{smallmatrix}0&M\\-M^t&0\end{smallmatrix}\right)$ と書く。
$M:\mathcal O(-2)^3\to\mathcal O(-1)^3$ の余核を $\theta$ とすると、
固定Gram表示の等方Grassmann閉包は
\[
 X_Q=\mathbb P_E(\mathcal O\oplus\theta)\cup\mathbb P^2_+\cup\mathbb P^2_-
\]
という三成分単純正規交差曲面である。
二重曲線は主成分の二つの切断で、法束は $\theta,\theta^{-1}$ となる。
主成分の符号商は $\mathbb P_E(\mathcal O\oplus\theta^2)$ である。
\item 一般の交代二形式の三つ組への一次摂動は種数二十K3曲面へ平滑化する。
全空間の十八普通二重点を二平面のイデアルのblow-upで解消すると、基底変換なしで
\[
 \operatorname{Bl}_9\mathbb P^2\ \cup\ 
 \mathbb P_E(\mathcal O\oplus\theta)\ \cup\
 \operatorname{Bl}_9\mathbb P^2
\]
というType II Kulikovモデルを得る。
有理端成分の反標準束が半豊富であることは $\theta$ の有限位数性と同値である。
位数 $m$ はHalphen指数に一致する。
幾何学的一般点、また複素数体上の非常に一般の点では、両端の反標準束はnefで半豊富でない。
\item 複素数体上で、極限の重み二格子は $U\oplus E_8(-1)^{\oplus2}$ であり、
Pl\"ucker類の平方は三十八である。
その直交格子の平方 $-2$ の根系は $A_8\oplus A_8$ であり、全格子を整数Gram行列と貼合せ指数で指定できる。
モノドロミーの対数は $N^2=0$、階数二、最大整数因子二を持つ。
\item 複素数体上で、二つの平滑化切断への写像は
\[
 (\Lambda^2(T\oplus B)^*)^{\oplus3}\longrightarrow
 H^0(E,H^3\theta)\oplus H^0(E,H^3\theta^{-1})
\]
へ全射する。原点と九点の標識を選び、
$\mathcal P_E=\ker(\Sigma:E^{18}\to E)$ と置くと、境界周期は
\[
 \mathcal P_E/\Delta E[3]\simeq\operatorname{Hom}(L/\mathbb Zh,E)
\]
と明示される。逆Gramパラメータと平滑化方向を動かす族は、この境界型の
十八次元の偏極付き周期空間を支配する。
固定したある楕円曲線上で、十分大きい全ての整数を両端のHalphen指数として実現できる。
\end{enumerate}
\end{theorem}

対称行列式と線束の対応はBeauvilleの理論を用いる\cite{BeauvilleDeterminantal}。
一般の交代二形式の三つ組に付随する楕円スクロールは古典的であり\cite{BazanMezzetti}、
交代二形式の一般の三つ組による種数二十K3のGrassmann表示はMukaiの構成である\cite{MukaiGenus20}。
Type IIの一般理論とその格子記述も既知である\cite{KulikovDegenerations,PerssonPinkhamDegenerations,KondoTypeII}。
本稿では、全自己反対方程式からこれらの対象への写像を作り、
特殊な直和形の交代二形式の三つ組の全閉包、二つの線織面の間の二乗商、行列式線束による端成分の判定、
この具体的退化の偏極付き整数格子、十八点配置による周期の整数係数式を計算する。
二つの完全線形系への全射性と次数九の同種写像を組み合わせ、この制約付き族が境界周期の全方向を与えることを示す。

Hermite曲線の同型類はnetに依存するが、次数零線束 $\theta$ はnetを固定しても変化する。
この支配性と、それに伴うモジュライの二重被覆は別稿\cite{PaperVI}で証明する。
同稿は古典的net写像 $\mathcal R_3\dashrightarrow Y_0(2)$ の相対有理性と、
粗モジュラー関数体上のBrauer障害・指数六も扱う。
本稿の幾何学的な射影直線束表示は、同関数体上での自明な束の存在を意味しない。
署名 $(3,3)$ の六次Prym復元、十九変数・九方程式、および標準的七次被覆は別稿\cite{PaperVII}にまとめた。

\section{全自己反対方程式と標識付き有理商}
\label{sec:prym-setup-v26}
基礎体は $\mathbb C$ とし、有限計算の入力は $\mathbb Q$ 上に取る。
$L(x)=\sum_{i=1}^3x_iL_i$ は六行七列の一次式行列である。
四頂点箙の矢の本数は $(3,6,3)$、左右の二次関係の次元はともに七であり、
右平面との同一視の後に右関係を $R_i=JL_i$ と書く。
\[
 J=\operatorname{diag}(I_p,-I_q),\qquad
 L=\begin{pmatrix}P\\B\end{pmatrix},\qquad A=A^t\in\operatorname{GL}_7.
\]
課す方程式は
\begin{equation}\label{eq:full-incidence-v26}
 L_iAL_j^tJ=JL_jAL_i^t\qquad(1\leq i,j\leq3)
\end{equation}
の\emph{全て}である。ブロックで書けば
\begin{align*}
 P_iAP_j^t&=P_jAP_i^t,& B_iAB_j^t&=B_jAB_i^t,\\
 P_iAB_j^t+P_jAB_i^t&=0 &&(i\leq j)
\end{align*}
となる。同符号の式を省いて混合直交性だけを使うと、係数 Gram 行列も逆構成も変わる。

対象の同値は、関係枠の $\operatorname{GL}_7$、二つの行ブロック内の
$\operatorname{GL}_p\times\operatorname{GL}_q$、共通の平面座標変更、
および pairing $A$ の非零スカラー倍で定める。
関係枠 $L\mapsto LM$ に対して $A\mapsto M^{-1}AM^{-t}$ とする。
有理商という語は、この作用の不変関数体を持つ幾何的な商の非空開部分を意味する。
普遍代数が粗モジュライ全体に降下するという主張は含めない。
$p=q=3$ では、まず $P$ に用いる固有空間を指定した商を使う。
六次の場合のブロックの順序と内在的復元は、別稿\cite{PaperVII}で扱う。

\begin{proposition}[既存成分と開条件]
\label{prop:prym-baseline-v26}
別稿\cite{PaperII}で構成した固定成分 $Z_1,Z_2,Z_3$ のうち、
$p=3,2$ に対応する成分を、それぞれ $Z_1,Z_3$ と書く。
それらは、次の条件を同時に満たす非空開部分を持つ。
\begin{enumerate}
\item $P(x)$ は平面の全点で階数 $p$ を持ち、
$Z=V(I_6(L))$ は相異なる二十一点からなる。
\item 線形核から右関係を復元する写像の階数は所要の最大値であり、
pairing はスカラー倍を除いて一意である。
従って上の枠付き表示は、指定反対合を保つ関係代数の分類に対応する。
\item 点対合の固定点は $B$ の階数低下点であり、固定点数は
$f(p)=21-12p+2p^2$、非固定点数は $2pq$ である。
非固定点は $C$ 上にあり、固定点は $C$ の外にある。
\end{enumerate}
対応する標識付き有理商の次元は
\[
 \begin{array}{c|c|c|c|c}
 \text{成分}&(p,q)&\text{有理商}&\text{次元}&(\text{固定点},\text{非固定点})\\\hline
 Z_3&(2,4)&\mathfrak Z_3&8&(5,16)\\
 Z_1&(3,3)&\mathfrak Z_1^+&7&(3,18)
 \end{array}
\]
である。$A_{77}=1$ の全incidenceチャートは、それぞれ
百五十三変数と六十九式、百五十三変数と七十二式で定義され、
境界点 $B_-$ で滑らかである。
\end{proposition}

\begin{proof}[用いる入力と幾何学的内容]
関係成分の構成、商の次元および同じ成分での被約点配置の存在は
\cite[反対合の固定部分族・全署名の点配置の各節]{PaperII}を用いる。
付属の \nolinkurl{hensel_other_signatures.py} と二つの保存証明書は、
その有限入力を完全に含む。
Hensel 分枝は、滑らかな境界点で自由変数を指定し、単元 Jacobian を持つ残りの変数を一意に解いたものである。
本稿の新しい計算は、基点、pivot、自由変数を全て照合した同一分枝上で行う。
単に同じ素数を法として合同な別成分の点を選んでいるのではない。

計数の内容を確認しておく。混合直交性により、$B^t$ は階数 $q+1$ の束
\[
 \mathcal E=\ker\bigl(\mathcal O(1)^7\xrightarrow{PA}\mathcal O(2)^p\bigr)
\]
の $q$ 本の切断を与える。
その有限な階数低下スキームの長さは
\[
 c_2(\mathcal E)=[H^2]\frac{(1+H)^7}{(1+2H)^p}=21-12p+2p^2
\]
である。非固定点では $P$ と $B$ の行空間が非零に交わり、
直交性から $P$ の Gram 形式が退化する。
固定点が同時にその退化曲線上にあれば、$I_6(L)$ の法線微分が階数二を持てず、
$Z$ の被約性に反する。これが曲線による二十一点の分離を与える。
\end{proof}


\section{滑らかなGram曲線と点配置に付随する二重被覆}
\label{sec:gram-prym-v25}

命題\ref{prop:prym-baseline-v26}の行列式曲線を、線束を伴う滑らかな曲線へ改善する。
全自己反対incidenceと指定反対合の枠を固定する。
対称行列表示と二捩れ線束の対応は既知の理論である
\cite[\S2, Proposition 4.6]{BeauvilleDeterminantal}。
ここで証明する本族に固有の点は、既存の固定成分内で滑らかさが実現することと、
その線束を二十一点配置の非固定部分から復元できることである。

\begin{theorem}[滑らかな四次・六次曲線の実在]
\label{thm:gram-smooth-v25}
$Z_3$ と $Z_1$ は、それぞれ、命題\ref{prop:prym-baseline-v26}の被約二十一点配置の条件に加えて
\[
 C=\Gamma_P=V(\det(PAP^t))
\]
が滑らかな平面四次曲線、滑らかな平面六次曲線となる非空開部分を持つ。
従ってこれらの曲線は幾何学的に既約であり、種数はそれぞれ三、十である。
$Z_2$ の非退化二次曲線と合わせ、三族とも滑らかなGram曲線を持つ開部分が存在する。
\end{theorem}

\begin{proof}
命題\ref{prop:prym-baseline-v26}と同じ素数 $101$、方向選択seed、基点、pivot変数、自由変数を用いる。
固定成分 $Z_3,Z_1$ のHensel解を精度 $101^{256}$ まで再構成する。
自由変数の全リストを既存の被約二十一点証明書と照合しているので、
これらは命題\ref{prop:prym-baseline-v26}で同一成分への所属と被約二十一点性を証明した、
同じ代数的 $101$ 進点の有限精度表示である。

$F=\det(PAP^t)$、$d=\deg F=2p$ と置く。
$F_x,F_y,F_z$ の全てに次数 $2d-4$ の単項式を掛けたものを、
次数 $3d-5$ の単項式基底で並べた行列を $\mathcal M_F$ とする。
次の小行列式を証明書で指定する。
\[
\begin{array}{c|c|c|c|c|c}
 \text{成分}&d&3d-5&\mathcal M_F\text{ の大きさ}&
 \text{指定小行列式の付値}&\text{保証精度}\\\hline
 Z_3&4&7&45\times36&11&256\\
 Z_1&6&13&135\times105&57&256
\end{array}
\]
前者は三十六次、後者は百五次の小行列式である。
$F$ の係数、偏微分、Macaulay行列、小行列式はいずれも
Hensel座標の整係数多項式であり、分母を導入しない。
従って入力の $101^{256}$ 精度が小行列式にもそのまま保証される。
付値が保証精度未満なので、真のHensel解で小行列式は非零である。

よって $F_x,F_y,F_z$ が生成する斉次イデアルは
次数 $3d-5$ の全ての単項式を含む。
もし代数閉包上の射影点で三偏微分が同時に零になれば、
その点では同次数の全ての単項式が零となり矛盾する。
従って射影的な共通零点はなく、$C$ は幾何学的に滑らかである。
正標数での退化した基点に滑らかさを主張しているのではなく、
標数零のHensel点における非消滅を証明していることに注意する。

滑らかさは開条件なので、同じ既存成分上に求める非空開部分を得る。
滑らかな正次数平面曲線は、異なる正次数成分があれば交点を持って特異となるため既約である。
種数は平面曲線の随伴公式から $(d-1)(d-2)/2$ である。
\end{proof}

\paragraph{再現入口.}
\nolinkurl{gram_curves_v25.py} と
\nolinkurl{gram_curves_v25_certificate.json} が指定小行列式の行番号、
正規化した行列式の剰余、Gram多項式の係数hash、既存証明書との同一性を記録する。
通常の実行は読取り専用である。
この検証は滑らかさのための追加の有限証明書であり、
同一成分への所属と被約二十一点性の証明は命題\ref{prop:prym-baseline-v26}のものを使用する。

\begin{proposition}[非自明二捩れ線束と非固定点の除数]
\label{prop:gram-prym-divisor-v25}
命題\ref{prop:prym-baseline-v26}の一般署名命題の仮定に加えて、$C=\Gamma_P$ が滑らかであるとする。
$p=1,2,3$、$q=6-p$ とし、$i:C\hookrightarrow\mathbb P^2$ と書く。
対称行列 $Q_P=PAP^t$ の余核を
\begin{equation}\label{eq:gram-line-resolution-v25}
 0\longrightarrow\mathcal O_{\mathbb P^2}(-2)^p
 \xrightarrow{Q_P}\mathcal O_{\mathbb P^2}^{p}
 \longrightarrow i_*\mathcal L\longrightarrow0
\end{equation}
で定め、$\eta=\mathcal L(1-p)$ と置く。
このとき $\mathcal L$ は $C$ 上の線束であり、
\[
 \eta^{\otimes2}\simeq\mathcal O_C.
\]
$p=2,3$ では $H^0(C,\eta)=0$ なので $\eta$ は非自明である。
$p=1$ では $\eta=\mathcal O_C$ となる。

非固定点を被約有効除数
\[
 D_{\rm mov}=\sum_{x\in Z\setminus\operatorname{Fix}(\iota)}x
\]
と見なすと、
\begin{equation}\label{eq:moving-prym-divisor-v25}
 \mathcal O_C(D_{\rm mov})\simeq\eta(q),
 \qquad
 \eta\simeq\mathcal O_C(D_{\rm mov}-qH)
\end{equation}
となる。従ってこの二捩れ線束は、指定反対合に付随するGram曲線と
非固定点の除数からも復元できる。
\end{proposition}

\begin{proof}
$C$ 上で $Q_P$ の階数は常に $p-1$ である。
実際、階数が $p-2$ 以下なら随伴行列が零で、
$\det Q_P$ の全微分が零となり $C$ の滑らかさに反する。
局所的に単元ブロックを消すと $Q_P$ は
$\operatorname{diag}(I_{p-1},f)$ と同値になり、$f$ は $C$ の局所方程式である。
従って余核は $C$ 上階数一の局所自由加群である。

式\eqref{eq:gram-line-resolution-v25}を
$\mathcal H om(-,\mathcal O_{\mathbb P^2}(-2))$ で双対化する。
$Q_P$ の対称性により、得られる余核も $i_*\mathcal L$ である。
一方、次数 $2p$ のCartier因子の双対性から
\[
 \mathcal E xt^1_{\mathbb P^2}
 (i_*\mathcal L,\mathcal O_{\mathbb P^2}(-2))
 \simeq i_*\bigl(\mathcal L^{-1}(2p-2)\bigr).
\]
よって $\mathcal L^{\otimes2}\simeq\mathcal O_C(2p-2)$、
従って $\eta^{\otimes2}\simeq\mathcal O_C$ である。
分解を $1-p$ だけねじると
\[
 0\longrightarrow\mathcal O(-p-1)^p
 \longrightarrow\mathcal O(1-p)^p
 \longrightarrow i_*\eta\longrightarrow0.
\]
$p=2,3$ では $H^0(\mathcal O(1-p))=0$ と
$H^1(\mathbb P^2,\mathcal O(-p-1))=0$ により $H^0(C,\eta)=0$ を得る。
$p=1$ は一行一列の表示そのものである。

次に $C$ 上の自明束 $\mathcal V=\mathcal O_C^7$ に $A$ の非退化対称形式を入れ、
$P$ の行空間を部分束
\[
 \mathcal S\simeq\mathcal O_C(-1)^p\subset\mathcal V
\]
とする。その根基 $\mathcal T=\mathcal S\cap\mathcal S^\perp$ は線束である。
式\eqref{eq:gram-line-resolution-v25}を $C$ に制限した際のTor項を計算し、
一つねじると
\[
 \mathcal T\simeq\mathcal L(1-2p)=\eta(-p)
\]
を得る。

命題\ref{prop:prym-baseline-v26}により $B$ の階数低下点は固定点に一致し、固定点は $C$ の外にある。
従って $C$ 上では $B^t$ の像も部分束
\[
 \mathcal B\simeq\mathcal O_C(-1)^q\subset\mathcal S^\perp
\]
となる。直交性は全incidenceの混合ブロックから従う。
商 $\mathcal F=\mathcal S^\perp/\mathcal B$ は線束であり、
\[
 \det(\mathcal S^\perp)\simeq\mathcal O_C(-p),\qquad
 \mathcal F\simeq\mathcal O_C(q-p)
\]
である。従って自然な射
\[
 \mathcal T\longrightarrow\mathcal F
\]
は $\eta^{-1}(q)\simeq\eta(q)$ の切断を定める。
その零点は $\mathcal T(x)\subset\mathcal B(x)$ となる点である。
これは $\mathcal S(x)\cap\mathcal B(x)\ne0$、すなわち
$\operatorname{rank}L(x)<6$ と同値であり、$C$ 上では非固定点ちょうど全部である。
点スキームが有限なので切断は恒等的には零でない。
この切断の零除数の次数は
\[
 \deg\eta(q)=2pq
\]
で、既に確定した相異なる非固定点数も $2pq$ である。
従って各零点の重複度は一であり、式\eqref{eq:moving-prym-divisor-v25}を得る。
\end{proof}

\begin{corollary}[Prym曲線とPrym-canonical除数]
\label{cor:gram-prym-canonical-v25}
上の非空開部分で、$Z_3$ から種数三の曲線と非自明二捩れ線束、
$Z_1$ から種数十の曲線と非自明二捩れ線束を得る。
対応する連結エタール二重被覆の種数はそれぞれ五、十九である。
特に $Z_1$ では
\[
 K_C\simeq\mathcal O_C(3),\qquad
 D_{\rm mov}\in|K_C\otimes\eta|,\qquad
 h^0(C,K_C\otimes\eta)=9.
\]
従って十八個の非固定点は、Prym-canonical線形系の被約除数をなす。
\end{corollary}

\begin{proof}
$\eta^{\otimes2}\simeq\mathcal O_C$ の同型を選ぶと
$\operatorname{Spec}_C(\mathcal O_C\oplus\eta)$ はエタール二重被覆となる。
$\eta$ が非自明なので連結であり、Riemann--Hurwitzから種数は $2g(C)-1$ である。
六次曲線では随伴公式により $K_C=\mathcal O_C(3)$ であり、$q=3$ なので
式\eqref{eq:moving-prym-divisor-v25}を適用できる。
Riemann--Rochと $h^0(C,\eta)=0$ により
$h^0(C,K_C\otimes\eta)=g(C)-1=9$ となる。
\end{proof}

\begin{proposition}[補助Gram曲線との二重交差]
\label{prop:complementary-gram-contact-v25}
上と同じ仮定の下で
\[
 \Gamma_B=V(\det(BAB^t))
\]
は次数 $2q$ の非零行列式曲線であり、$C$ を成分に持たない。
$C$ 上のCartier除数として
\begin{equation}\label{eq:complementary-gram-contact-v25}
 \Gamma_B|_C=2D_{\rm mov}
\end{equation}
が成立する。特に $Z_1$ では二つの六次行列式曲線の交差は、
十八個の非固定点に重複度二ずつ集中する。
また、全ての固定点は $\Gamma_B$ の特異点である。
補助曲線の既約性と特異点は第\ref{sec:complementary-nodes-v26}節で決定する。
\end{proposition}

\begin{proof}
$\mathcal S^\perp$ の根基も $\mathcal T$ なので、
\[
 \mathcal H=\mathcal S^\perp/\mathcal T
\]
は階数 $q$ の非退化対称束であり、$\det\mathcal H\simeq\eta^{-1}$ である。
$\mathcal B\to\mathcal H$ は同階数間の射で、その行列式は
$\eta^{-1}(q)$ の切断である。
これが退化するのは $\mathcal T\cap\mathcal B\ne0$ の点なので、
前の命題と同じ次数計算により、その零除数は $D_{\rm mov}$ である。
$\mathcal B$ 上のGram形式は、$\mathcal H$ の非退化形式をこの射で引き戻したものである。
従ってその行列式の零除数は二倍され、
式\eqref{eq:complementary-gram-contact-v25}を得る。
非固定点は有限なので行列式の $C$ への制限は恒等的には零でない。
これにより多項式自体が非零であり、$C$ を成分に持たないことも分かる。

固定点では $B$ の階数は $q-1$ である。
$G=BAB^t$ の階数が $q-2$ 以下なら、$\operatorname{adj}(G)=0$ なので
$d(\det G)=0$ である。
階数が $q-1$ なら、その核の生成元 $v$ は $B^tv=0$ を満たし、
\[
 v^t(dG)v=v^t(dB)AB^tv+v^tBA(dB)^tv=0.
\]
随伴行列は $vv^t$ の非零倍なので、この場合にも $d(\det G)=0$ となる。
従って固定点は $\Gamma_B$ の特異点である。
\end{proof}

\begin{remark}[二つの二重構造の区別]
$\eta$ が定めるエタール二重被覆は、曲線 $C$ の上の被覆である。
一方、指定反対合は $D_{\rm mov}$ の点を二点ずつ交換する。
この二つを同一視する主張はしていない。
$p=3$ の対称 $3\times3$ 二次式行列 $Q_P$ は、
$\mathbb P^2\times\mathbb P^2$ の双次数 $(2,2)$ 超曲面
$y^tQ_P(x)y=0$ も定め、その第一射影のconic bundleの判別式は $C$ である。
そのPrymデータは上の $\eta$ で記述される既知の構成に属する
\cite[\S1--2]{BoehningBothmerPrym}。
本族が一般の双次数 $(2,2)$ 超曲面を支配するとは主張せず、
一般のconic bundleの非安定有理性をこの特別な部分族へ転用しない。
\end{remark}

\section{種数三Prymモジュライへの支配性}
\label{sec:prym-dominance-v26}

本節では $Z_3$ の署名を $(p,q)=(2,4)$ として、
指定反対合を保つ同型のもとで商を取る。
この有理商を $\mathfrak Z_3$ と書く。
命題\ref{prop:prym-baseline-v26}より
$\dim\mathfrak Z_3=35-27=8$ である。
滑らかな種数三曲線と非自明二捩れ線束の対 $(C,\eta)$ のモジュライを
$\mathcal R_3$、その非超楕円部分を $\mathcal R_3^{\rm nh}$ とする。

\begin{theorem}[Prymモジュライへの支配性]
\label{thm:prym-dominance-v26}
第\ref{sec:gram-prym-v25}節の構成は支配的な有理写像
\[
 \Phi:\mathfrak Z_3\dashrightarrow\mathcal R_3
\]
を定める。その像の次元は六であり、一般ファイバーは二次元である。
この写像の定義域には、二十一点が被約で、Gram四次曲線が滑らかである非空開部分を取れる。
主張は指定反対合を持つ対象のモジュライについてのものである。
\end{theorem}

証明の中心は、Prymモジュライへ移る前のGram行列そのものが一般に動くという、
次のより強い局所計算である。

\begin{proposition}[十八個のGram係数の独立性]
\label{prop:gram-differential-v26}
全incidenceの正規化された枠付きチャートで
\[
 \gamma:(L,A)\longmapsto Q_P=PAP^t
 \quad\in\quad
 \operatorname{Sym}^2(k^2)\otimes H^0(\mathbb P^2,\mathcal O(2))
 \simeq\mathbb A^{18}
\]
を考える。
既存の $Z_3$ のHensel点で $d\gamma$ は階数十八を持つ。
従ってこの既約成分から対称二行二列二次式行列の空間への写像は支配的である。
\end{proposition}

\begin{proof}
変数は $L$ の百二十六係数と、$A_{77}=1$ と正規化した
対称 $7\times7$ 行列の二十七係数、合計百五十三個である。
全incidenceは六十九式 $F=0$ で定義される。
命題\ref{prop:prym-baseline-v26}の素数 $101$、seed $17029$ のHensel点を用い、
基点・pivot変数・自由変数を既存の証明書と全て照合する。
六十九個のpivot変数を $z$、八十四個の自由変数を $t$ とすると、
基点での $J_z=\partial F/\partial z$ の行列式は $12288$ であり、
$101$ 進単元である。
従って真のHensel点の接空間は
\[
 dz=-J_z^{-1}J_t\,dt
\]
で与えられる。
$Q_P$ の十八係数の変数全体に関する微分を $D=(D_z,D_t)$ と置けば、
制限された微分は
\begin{equation}\label{eq:gram-differential-schur-v26}
 G=D_t-D_zJ_z^{-1}J_t\in\operatorname{Mat}_{18\times84}(\mathbb Z_{101})
\end{equation}
となる。

精度 $101^{32}$ でHensel点を再構成した後、
最後に補正された座標で $J_z^{-1}$ を更新し直し、
$J_zJ_z^{-1}=I$ を同精度で検証する。
式\eqref{eq:gram-differential-schur-v26}の指定十八次小行列式 $\Delta$ は
\[
 v_{101}(\Delta)=2<32,
 \qquad 101^{-2}\Delta\equiv914960\pmod{101^3}
\]
を満たす。
選択した列は有限証明書に明記されており、消去のpivot付値は
$0$ が十七個と $2$ が一個である。

ここで $J_z^{-1}$ は $\mathbb Z_{101}$ 上の行列であり、
式\eqref{eq:gram-differential-schur-v26}は精度を失わない。
小行列式もその成分の整係数多項式なので、
真の小行列式は計算値と $101^{32}$ を法として一致する。
従って有限精度で得た付値二は真の付値であり、$\Delta\ne0$ である。
よって微分の階数は十八である。
標数零の滑らかな点での微分の全射性から、この成分上の $\gamma$ は支配的である。

この計算では、係数多項式の微分を、独立に組み立てた行列積
$PAP^t$ の方向微分とも照合している。
行列積は各方向について三次以下なので、四つの整数点での補間が
微分の正確な照合を与える。
\end{proof}

\begin{lemma}[四次曲線と二捩れ線束の二次式表示]
\label{lem:quartic-prym-representation-v26}
標数零の代数閉体上で、滑らかな非超楕円種数三曲線 $C$ と
$0\ne\eta\in\operatorname{Pic}^0(C)[2]$ の任意の対は、
対称 $2\times2$ 二次式行列の行列式とその余核から得られる。
\end{lemma}

\begin{proof}
これは対称行列表示の一般論の特別な場合である
\cite[\S2.3, Proposition 4.6]{BeauvilleDeterminantal}。
ここでは種数三の場合の直接的な構成を与える。
標準埋込み $C\hookrightarrow\mathbb P^2$ を取り、
$M=K_C\otimes\eta$ と置くとRiemann--Rochから $h^0(C,M)=2$ である。

このpencilは基点を持たない。
実際、$a$ が基点ならRiemann--Rochより
$h^0(C,\eta(a))=1$ であり、ある $b\in C$ に対して
$\eta\simeq\mathcal O_C(b-a)$ となる。
$a=b$ は $\eta$ の非自明性に反し、$a\ne b$ なら
$2a\sim2b$ が次数二の写像 $C\to\mathbb P^1$ を与え、非超楕円性に反する。

基底 $s_0,s_1$ と $\eta^{\otimes2}\simeq\mathcal O_C$ を選ぶ。
制限写像
\[
 H^0(\mathbb P^2,\mathcal O(2))\xrightarrow{\sim}H^0(C,K_C^{\otimes2})
\]
は同型なので、$s_0^2,s_0s_1,s_1^2$ は一意に二次式
$a,b,c$ に持ち上がる。
$ac-b^2$ は $C$ 上で零である。
これが零多項式なら、多項式環の一意分解性から
\[
 a=h\ell_0^2,\quad b=h\ell_0\ell_1,\quad c=h\ell_1^2
\]
と書ける。
次数から、$\ell_0,\ell_1$ が定数の場合と、$h$ が定数で
$\ell_0,\ell_1$ が一次式の場合だけがある。
前者は $s_0,s_1$ の比例を意味し、後者は
$s_0/\ell_0$ が $\eta$ の零点も極もない切断であることを意味する。
いずれも矛盾する。
従って $ac-b^2$ は非零であり、$C$ の四次方程式の定数倍である。

$Q=\left(\begin{smallmatrix}a&b\\b&c\end{smallmatrix}\right)$ と置く。
$\mathcal O_C^2\to M$ を基底切断 $(-s_1,s_0)$ で定めると、
$Q$ の像はこの写像の核に含まれる。
基点がないため写像は全射であり、$Q$ は $C$ 上で階数一なので、
$\operatorname{coker}Q\simeq M$ となる。
従って余核を $\mathcal O_C(-1)$ でねじって得る線束は $\eta$ である。
\end{proof}

\begin{proof}[定理\ref{thm:prym-dominance-v26}の証明]
命題\ref{prop:gram-differential-v26}と
補題\ref{lem:quartic-prym-representation-v26}により、
枠付きincidence成分から $\mathcal R_3^{\rm nh}$ への写像は支配的である。
既存の有限被約二十一点条件とGram曲線の滑らかさは同じ成分の非空開条件であり、
それらに制限しても支配性は保たれる。

枠を外す際のスカラーを確認する。
関係基底変換は $P$ と $A$ を逆向きに変換するため、
$A_{77}=1$ の再正規化を除けば $Q_P$ を保つ。
再正規化による非零スカラー倍は、行列式曲線と余核線束を変えない。
正固有空間の $\operatorname{GL}_2$ は $Q_P$ に合同変換として作用し、
負固有空間の $\operatorname{GL}_4$ は $Q_P$ を変えず、
外側の $\operatorname{PGL}_3$ は標準平面の座標を変える。
従って $(C,\eta)$ の同型類は指定反対合を保つ基底変換に不変であり、
$\Phi$ は $\mathfrak Z_3$ 上に降下する。

非超楕円種数三曲線のモジュライは六次元で、非自明二捩れ線束の選択は有限である。
よって $\dim\mathcal R_3=6$ である。
$\dim\mathfrak Z_3=8$ と支配性により、一般ファイバーの次元は $8-6=2$ となる。
\end{proof}

\paragraph{有限証明書とファイバーの局所計算.}
\nolinkurl{gram_dominance_v26.py} と
\nolinkurl{gram_dominance_v26_p2_certificate.json} は、上の微分に加えて、
後のファイバー解析に使う次の有限計算を記録する。
$P$ の六つの係数ベクトルの張る空間を $H\subset k^7$ とし、
\[
 K=\{b(x)\in k^7\otimes H^0(\mathbb P^2,\mathcal O(1)):
                      P(x)A b(x)^t=0\}
\]
と置く。
同じHensel点で、以下が成り立つ。
\[
\begin{array}{c|c|c}
 \text{対象}&\text{指定小行列式の付値}&\text{帰結}\\\hline
 P\text{ の }6\times7\text{ 係数行列}&0&\dim H=6\\
 A|_H\text{ の }6\times6\text{ Gram行列}&1&A|_H\text{ は非退化}\\
 PAb^t=0\text{ の }12\times21\text{ 係数行列}&0&\dim K=9\\
 B\text{ が定める }4\times9\text{ 行列}&0&W=\langle B\rangle\in\operatorname{Gr}(4,K)\\
 \text{等方性方程式の }18\times20\text{ 微分}&2&\text{零点スキームは当該点で滑らかな曲面}
\end{array}
\]
最後の十八式は、三つの座標対 $0\le i<j\le2$ に対して
\[
 \omega_{ij}(b,c)=b_iAc_j^t-b_jAc_i^t
\]
という $K$ 上の交代形式を取り、$\omega_{ij}|_W=0$ としたものである。
その指定小行列式の正規化した剰余は
$101^{-2}\Delta_{\rm Gr}\equiv102175\pmod{101^3}$ である。
使う三つの基底選択の小行列式は全て単元なので、この計算にも精度損失はない。
これは当該開部分の局所的な滑らかさを証明する。
Grassmann多様体内の零点スキーム全体の滑らかさや、
その射影的閉包の幾何は、この有限証明書だけからは結論しない。

さらに、$n\in H^\perp$ を $P$ の係数行列と $A$ の積の六次小行列式で取ると、
四行三列の行列 $(B_{r,i}An)$ は階数三である。
行 $0,1,3$ の小行列式の付値は一で、正規化剰余は
$555564\pmod{101^3}$ である。
これは $W\to H^\perp\otimes H^0(\mathcal O(1))$ の像が三次元であることを保証する。

% Standalone fragment, to be integrated by the root agent.
\section{四次Gram曲線を固定したファイバーと補助楕円曲線}
\label{sec:prym-fiber-v26}
\label{subsec:prym-elliptic-ruled-v26}

三つの交代形式から六次楕円スクロールを作ること自体は古典的であり、\cite{BazanMezzetti} で詳述されている。
本節の目的は、その既知のスクロールと、指定反対合を保持した本族のPrymファイバーとの双有理同定である。
以下では標数零の代数閉体上で議論する。
$Q=PAP^t$ は $2\times2$ 対称二次行列、$P$ は $2\times7$、$B$ は $4\times7$ とする。
$P$ の六本の係数ベクトルのGram行列を $G$ と書く。
同符号ブロックを含む全incidenceを用いることが本節でも不可欠である。

\begin{proposition}[固定GramファイバーのGrassmann表示]
\label{prop:gram-fiber-grassmann-v26}
$\det G\ne0$ の開部分では、関係枠を変えることにより
\[
 P(x)=\begin{pmatrix}x_1&x_2&x_3&0&0&0&0\\
 0&0&0&x_1&x_2&x_3&0\end{pmatrix},\qquad
 A=\operatorname{diag}(G,1)
\]
とできる。$G$ の $3\times3$ ブロックを $G_{ab}$ と書くと、
$G_{11},G_{12}=G_{21},G_{22}$ はすべて対称であり、
$Q_{ab}(x)=x^tG_{ab}x$ によって $Q$ と $G$ は互いを一意に定める。

$[x]_\times c=x\times c$ とし、
\[
 D_i=\operatorname{diag}([e_i]_\times,[e_i]_\times),\qquad
 \Omega_{ij}=D_i^tG^{-1}D_j-D_j^tG^{-1}D_i\quad(i\ne j)
\]
と置く。$K=K_6\oplus K_3=k^6\oplus k^3$ 上に三つの交代形式
\[
 \omega_{ij}((c,\ell),(c',\ell'))
 =c^t\Omega_{ij}c'+\ell_i\ell'_j-\ell_j\ell'_i
\]
を定める。このとき、$Q$ を固定した全incidenceの解を関係枠と $B$ の行基底で割った空間は、
開条件を付した
\[
 X_Q/\langle\tau\rangle,\qquad
 X_Q=\{W\in\operatorname{Gr}(4,K):\omega_{ij}|_W=0\ (i<j)\},\qquad
 \tau(c,\ell)=(c,-\ell)
\]
で記述される。
\end{proposition}

\begin{proof}
係数六本の一次独立性から $P$ を表示の標準形にできる。
$P$--$P$ incidenceは $G_{12}$ が対称であることと同値である。
$G$ が可逆なので、$P$ の標準形を保つ関係枠変更により $A$ の右上六成分を消せる。
残るSchur補行列は非零の一成分であり、平方根を選んで一にする。
この標準化の残余は最後の座標の符号変更である。

一行の一次式 $b(x)=(b_6(x),\ell(x))$ に対する混合incidenceは
\[
 P(x)Ab(x)^t=0
\]
である。これは $G b_6(x)^t$ の二つの三成分ブロックが $x$ と直交することを意味し、
各ブロックは一意に $[x]_\times c_a$ と書ける。
従ってすべての解は
\[
 b(x)^t=\bigl(G^{-1}\operatorname{diag}([x]_\times,[x]_\times)c,\ \ell^tx\bigr),
 \qquad(c,\ell)\in K_6\oplus K_3
\]
となる。$B$ の四行は $K$ 内の四次元部分空間 $W$ を与える。
残る $B$--$B$ incidenceを係数で書けば、ちょうど三つの $\omega_{ij}|_W=0$ である。
最後の座標の符号変更が $\tau$ となる。
\end{proof}

\paragraph{正規化と次元.}
ここではpairing $A$ 自身を用いているので、$A_{77}=1$ を課す前のincidenceの次元は八十五である。
固定Gramファイバーでは
$85-18-49-16=2$ という次元勘定になる。
既存の $A_{77}=1$ チャートの次元八十四を用いる場合、
Gramを固定したまま動ける関係枠の次元は四十八である。
この一つのスカラーを二重に差し引かない。

\begin{proposition}[補助Pfaffian三次曲線の滑らかさ]
\label{prop:auxiliary-pfaffian-cubic-v26}
上記十八次元の $Q$ の空間の非空開部分で
\[
 E_Q=V\bigl(\operatorname{Pf}(z_1\Omega_{23}+z_2\Omega_{31}+z_3\Omega_{12})\bigr)
 \subset\mathbb P^2_z
\]
は滑らかな平面三次曲線である。
従って $E_Q$ 上の交代束写像
\[
 \Omega(z):\mathcal O_{E_Q}^{6}\longrightarrow\mathcal O_{E_Q}(1)^{6}
\]
は恒等的に階数四を持ち、その核 $\mathcal K_Q$ は階数二の束である。
\end{proposition}

\begin{proof}
非空性は一変数の族で直接証明できる。
$G(t)=\left(\begin{smallmatrix}A&C\\ C&D(t)\end{smallmatrix}\right)$ と置き、
\[
 A=\begin{pmatrix}0&1&0\\1&0&0\\0&0&1\end{pmatrix},\qquad
 C=\begin{pmatrix}1&0&0\\0&0&2\\0&2&0\end{pmatrix},\qquad
 D(t)=\begin{pmatrix}0&0&1\\0&t&0\\1&0&0\end{pmatrix}
\]
とする。各ブロックは対称であり、$\det G(t)=t-4$ である。
上の公式で交代形式を作ると、有理関数体上の恒等式として
\[
 \operatorname{Pf}\Omega(z)
 =-\frac{2tz_1^3+z_2^3+tz_3^3+(t-8)z_1z_2z_3}{t-4}
\]
を得る。$t=1$ では
\[
 E_Q:\quad 2z_1^3+z_2^3+z_3^3-7z_1z_2z_3=0
\]
である。
一般に $az_1^3+bz_2^3+cz_3^3+dz_1z_2z_3$ に対し $abc\ne0$ とすると、
特異点である射影点は三座標がすべて非零でなければならない。
三つの偏微分方程式を掛け合わせると $d^3+27abc=0$ が必要になる。
ここでは $(-7)^3+27\cdot2=-289\ne0$ なので、三次曲線は滑らかである。
滑らかさは開条件であり、求める非空開部分を得る。
Pfaffian恒等式は付属の短い有理関数体上の検証コードで再現できる。

六次交代行列のPfaffianの微分は四次Pfaffianの線形結合である。
従って $E_Q$ 上に階数二以下の点があれば全微分が零となり、滑らかさに矛盾する。
階数は四で一定であり、核は階数二の局所自由束になる。
\end{proof}

\begin{theorem}[固定Gramファイバーの楕円線織面モデル]
\label{thm:gram-fiber-elliptic-ruled-v26}
$E_Q$ が滑らかであるとする。
$X_Q$ の開部分で、次の三条件を課す。
\begin{enumerate}
\item $I_6(L)$ の零点集合が有限である。
\item 射影 $W\to K_3$ が階数三を持つ。
\item この開部分が非空である。
\end{enumerate}
このとき、この開部分を $\tau$ で割った空間は
\[
 S_Q=\mathbb P_{E_Q}(\mathcal K_Q)
\]
の非空開部分と双有理同値である。
特にその関数体は種数一曲線上の一変数有理関数体であり、幾何学的に既約な線織曲面である。
$S_Q$ は楕円曲線上の射影直線束である。

元の $Z_3$ 成分での開条件の非空性と、Prym 有理商の一般幾何学的ファイバーへの適用は、
系\ref{cor:coarse-prym-fiber-v26}で確認する。
\end{theorem}

\begin{proof}
まず $W\cap K_3=0$ を示す。
非零 $(0,\ell)\in W$ があれば $B$ の行基底を変えて、その一行を $\ell(x)e_7$ にできる。
直線 $\ell(x)=0$ 上ではその行が恒等的に零となるので、$L$ の階数は五以下となる。
従って $I_6(L)$ の零点集合が直線を含み、有限性に反する。

$W\to K_3$ は階数三なので、$k=W\cap K_6$ は一次元である。
$K_6$ への射影は $W$ 上単射であり、その像を $T$ とすると $\dim T=4$ である。
$\omega_{ij}|_W=0$ から
\[
 T\subset F_k:=\{v\in K_6:k^t\Omega_{ij}v=0\text{ for all }i<j\}.
\]
従って写像 $K_6\to k^3$, $v\mapsto(k^t\Omega_{23}v,k^t\Omega_{31}v,k^t\Omega_{12}v)$ の階数は高々二である。
ある $[z]\in\mathbb P^2$ に対して $k\subset\ker\Omega(z)$ となり、必然的に $[z]\in E_Q$ である。
この $[z]$ は一意である。二つあれば、その間の射影直線のすべてで $\Omega(z)$ は特異となり、
滑らかな三次曲線 $E_Q$ が直線を含むことになってしまう。
従って上記写像の階数はちょうど二であり、$\dim F_k=4$、$T=F_k$ となる。

以上により、$W$ は $S_Q$ 上の点 $([z],k)$ を定める。
逆に、この点に対して $F_k$ を作る。
各 $\Omega_{ij}|_{F_k}$ は $k$ を根基に含むので、
三次元空間 $U_k=F_k/k$ 上の三つの交代形式を与える。
これらが $\bigwedge^2U_k^*$ の基底をなすことは開条件である。
この条件は元の $W$ の存在から非空である。
実際、その $W$ は $F_k$ から $K_3$ への写像 $\phi$ のグラフであり、
$\ker\phi=k$、かつ $\phi^*\omega_{ij}^{(3)}=-\Omega_{ij}|_{F_k}$ を満たす。

従って逆構成は、与えられた可逆写像
$\bigwedge^2U_k\to\bigwedge^2K_3$ を $\bigwedge^2\bar\phi$ に持ち上げる問題になる。
三次元では
\[
 \operatorname{GL}(U_k,K_3)\longrightarrow
 \operatorname{GL}(\bigwedge^2U_k,\bigwedge^2K_3),\qquad
 \bar\phi\longmapsto\bigwedge^2\bar\phi
\]
は核 $\{\pm1\}$ の全射である。
基底を選べば行列式と逆転置による公式
$\bigwedge^2M=(\det M)M^{-t}$ がこれを示す。
よって持上げはちょうど二つ、互いに符号のみ異なる。
得られる二つのグラフは $\tau$ で交換される。
有限点配置その他の既存の開条件を戻しても、非空な部分は双有理型を変えない。
従って $\tau$ 商が $S_Q$ と双有理同値になる。
\end{proof}

\begin{remark}[得られた結論の範囲]
この定理は一般ファイバーの有理性を主張しない。
むしろ滑らかな射影モデルの不正則数は一であり、有理曲面にはならない。
また、$X_Q$ 全体が滑らかなK3曲面であるとも主張しない。
三つの一般の交代形式に対する同じGrassmann構成は種数二十のK3曲面を与えるが\cite{MukaiGenus20}、
本節の形式は $K_6\oplus K_3$ に関して特殊な直和形を持つ。
第\ref{sec:grassmann-boundary-v28}節で、この場合の閉包が三成分の単純正規交差曲面であることを示す。
従って一般の交代形式に関する滑らかさの結論を、そのまま本族へ適用することはできない。
\end{remark}

\paragraph{再現入口.}
\nolinkurl{verify_hesse_family.py} は $\mathbb Q(t)$ 上で三つの交代形式を作り、
上のPfaffian恒等式と $t=1$ における滑らかさの証明を再検証する。
\nolinkurl{hesse_family_certificate.json} に式と係数を記録する。
独立の有理係数例に対するMacaulay行列の階数検証は
\nolinkurl{verify_pfaffian_cubic.py} と \nolinkurl{pfaffian_cubic_certificate.json} に保存する。
いずれも通常実行は証明書との読取り専用照合である。

\subsection{Prym 有理商の一般幾何学的ファイバー}
\label{subsec:prym-descent-v26}

\begin{corollary}[粗ファイバーへの降下]
\label{cor:coarse-prym-fiber-v26}
写像 $\Phi:\mathfrak Z_3\dashrightarrow\mathcal R_3$ の一般幾何学的ファイバーは、
定理\ref{thm:gram-fiber-elliptic-ruled-v26}の
$S_Q=\mathbb P_{E_Q}(\mathcal K_Q)$ と双有理である。
従って幾何学的に既約であり、その滑らかな射影モデルは不正則数一を持つ線織曲面である。
\end{corollary}

\begin{proof}
まず必要な開条件が同じ成分で交わることを確認する。
命題\ref{prop:gram-differential-v26}により Gram 写像は十八次元の行列空間を支配する。
補助三次曲線の滑らかさは、その行列空間の非空開条件である。
一方、$\det G\ne0$、$W\to K_3$ の階数三、全二十一点の有限被約性は、
同じ既約成分に属する Hensel 点で確認された非空開条件である。
従ってこれらの交わりも非空であり、行列空間の稠密な開部分の上で
定理\ref{thm:gram-fiber-elliptic-ruled-v26}を適用できる。
同定理の逆構成では $S_Q$ の非空開部分が得られるので、
固定 Gram 行列の符号商の一般ファイバーは幾何学的に既約となる。

次に、Gram 行列の枠を外す際に余分な有限商が生じないことを確認する。
対称最小分解の一意性\cite[\S2.3]{BeauvilleDeterminantal}により、
同じ標準平面内の $(C,\eta)$ を表す二次対称行列は合同変換で結ばれる。
スカラーを含めて扱うには $[Q]$ を射影化し、
$\operatorname{PGL}_2\times\operatorname{PGL}_3$ の作用を使えばよい。
アフィン係数空間上で $\operatorname{PGL}_3$ の線形化を暗黙に選ぶ必要はない。

一般の平面四次曲線は自己同型を持たない
\cite[\S3.1, Theorem 3.1 and Figure 1]{QuarticAutomorphisms}。
その開部分に制限すれば、平面座標に由来する有限安定化群はない。
さらに一般の $Q$ の三つの係数行列は $\operatorname{Sym}^2(k^2)$ を張るので、
$Q$ を固定する正ブロックの合同変換は $\pm I_2$ だけである。
$P\mapsto-P$ は、関係枠の $-I_7$ と $B$ の行基底変更 $-I_4$ で打ち消せるため、
この残余作用も $X_Q/\langle\tau\rangle$ 上では自明である。
従って Prym 有理商の一般幾何学的ファイバーに、これ以上の有限商は現れない。

$S_Q$ は滑らかな種数一曲線上の射影直線束なので、射影束の構造写像 $\pi$ に対して
$\pi_*\mathcal O_{S_Q}=\mathcal O_{E_Q}$、$R^1\pi_*\mathcal O_{S_Q}=0$ である。
従って $h^1(S_Q,\mathcal O_{S_Q})=h^1(E_Q,\mathcal O_{E_Q})=1$ となる。
不正則数は滑らかな射影曲面の双有理不変量なので、一般ファイバーは有理曲面ではない。
\end{proof}

\paragraph{補助曲線の選択に依存しない意味.}
係数 Gram 行列 $G$ は $Q$ の別々の偏極から自然に得られる。
$K_6$ は $P$ の係数空間に値を持つ一次の直交 syzygy 全体であり、
三つの交代形式は二つの syzygy の pairing を平面座標について反対称化したものである。
従って正ブロックと平面の基底変更は、$K_6$ と形式の空間の基底変更を誘導する。
Pfaffian 曲線は射影同型まで、核束の射影化も同型まで不変である。
この意味で一般の Prym 対 $(C,\eta)$ に補助種数一曲線 $E_Q$ が付随する。
特定の原点を選んだ楕円曲線の群構造を標準的に指定しているわけではない。

\paragraph{一般関数体上の降下.}
以上は一般幾何学的ファイバーの同定である。
第\ref{sec:hermite-net-v27}節では補助曲線をHermite曲線として内在的に記述し、
別稿\cite{PaperVI}では粗モジュライの関数体上の降下障害を決定する。
その障害は非零であり、幾何学的な射影直線束は一般関数体上で非分裂二次曲線束となる。

\section{補助楕円曲線のHermite表示と二次同種写像}
\label{sec:hermite-net-v27}

四次の場合の補助曲線には、二次曲線のnetから直接得られる意味がある。
本節では、逆Gram行列を含むPfaffianを、netの直線への制限の行列式と同一視する。
これにより、一般Prymファイバーの底曲線は、退化二次曲線の二つの直線成分を
選ぶ古典的なHermite曲線であることが分かる。
この曲線の古典的な定義、および一般netを三次曲線の極二次曲線系として
表す場合のCayleyanとの関係は、\cite[著者公開原稿 \S3.2.2, Exercise 3.12]{DolgachevCAG}を参照されたい。

\subsection{Pfaffianと制限行列式の恒等式}

\[
 Q=\begin{pmatrix}a&b\\b&c\end{pmatrix},\qquad
 a=x^tA_0x,\quad b=x^tA_1x,\quad c=x^tA_2x,
 \qquad G=\begin{pmatrix}A_0&A_1\\A_1&A_2\end{pmatrix}
\]
と書く。ここで $U$ は平面の基礎となる三次元ベクトル空間である。
$N=\langle a,b,c\rangle\subset\operatorname{Sym}^2U^*$ を三次元とする。
双対平面の点 $[z]$ は直線 $\ell_z=\mathbb P(\ker z)$ を表す。
双対平面上の普遍階数二部分束を $\mathcal S$ とすると、制限写像
\[
 r_N:N\otimes\mathcal O\longrightarrow\operatorname{Sym}^2\mathcal S^*
\]
の行列式は三次式である。基底と体積形式を選んだ表示を $H_N(z)$ と書く。
具体的には $z_3\ne0$ で
\[
 x=(z_3u,z_3v,-z_1u-z_2v)
\]
を代入し、$a,b,c$ の係数を $u^2,uv,v^2$ の順に並べた三行三列行列を $R_N(z)$ とすると、
\[
 H_N(z)=\frac{\det R_N(z)}{z_3^3}.
\]
これは分母を持たない斉次三次式である。
$H_N(z)=0$ は、netの非零な二次曲線が直線 $\ell_z$ を含むことと同値である。

\begin{proposition}[Hermite--Pfaffian恒等式]
\label{prop:hermite-pfaffian-v27}
$\det G\ne0$ とする。第\ref{sec:prym-fiber-v26}節と同じ規約で
$\Omega(z)=z_1\Omega_{23}+z_2\Omega_{31}+z_3\Omega_{12}$ を作ると、
\[
 \boxed{\quad 2\det(G)\operatorname{Pf}\Omega(z)=H_N(z).\quad}
\]
従って補助曲線 $E_Q$ はnet $N$ のHermite曲線と、双対平面内で一致する。
特に $E_Q$ は $Q$ の三つの成分の選んだ基底よりも、net $N$ 自体に依存する。
\end{proposition}

\begin{proof}
まず $z=(0,0,1)$ とする。$M=G^{-1}$、行列の添字は一から六までとする。
Pfaffianを展開すると
\[
 \operatorname{Pf}\Omega_{12}
 =\det M_{(1,3,6),(5,3,6)}
  -\det M_{(2,3,6),(4,3,6)}.
\]
Jacobiの相補小行列式公式から
\begin{align*}
 \det(G)\operatorname{Pf}\Omega_{12}
 &=\det G_{(1,2,5),(1,4,5)}
   -\det G_{(1,2,4),(2,4,5)}\\
 &=\det\begin{pmatrix}
 (A_0)_{11}&(A_1)_{11}&(A_2)_{11}\\
 (A_0)_{12}&(A_1)_{12}&(A_2)_{12}\\
 (A_0)_{22}&(A_1)_{22}&(A_2)_{22}
 \end{pmatrix}.
\end{align*}
最後の等号には $A_1$ の対称性を用いた。
直線 $x_3=0$ への制限行列では中央の行の係数が二倍となるので、
これは $H_N(0,0,1)/2$ である。

任意の座標変更 $x=gx'$ に対し、$N^g$ と $G^g$ を変更後の対象とすると
\[
 H_{N^g}(z)=\det(g)^3H_N(g^{-t}z),\qquad
 \det G^g\operatorname{Pf}\Omega^g(z)
 =\det(g)^3\det G\operatorname{Pf}\Omega(g^{-t}z).
\]
第一式は制限束の行列式から従う。第二式は
$[gv]_\times=\det(g)g^{-t}[v]_\times g^{-1}$ と
$\operatorname{Pf}(T^t\Omega T)=\det(T)\operatorname{Pf}(\Omega)$ から従う。
従って一つの直線での恒等式を座標変更すれば全ての直線で成立し、求める三次式の恒等式を得る。
\end{proof}

\subsection{判別式三次曲線と非自明二捩れ}

netのパラメータ平面内に
\[
 \Delta_N=V\bigl(\det(\lambda_0A_0+\lambda_1A_1+\lambda_2A_2)\bigr)
 \subset\mathbb P(N)
\]
を取る。これは退化二次曲線をパラメータ化する判別式三次曲線である。
$E_Q$ と $\Delta_N$ の存在する平面は異なることに注意する。

\begin{theorem}[Hermite曲線のエタール二重被覆]
\label{thm:hermite-double-cover-v27}
$E_Q$ と $\Delta_N$ が滑らかな開部分では、直線成分をそれを含む退化二次曲線へ送る写像
\[
 \pi_N:E_Q\longrightarrow\Delta_N
\]
は連結な有限エタール二重被覆である。
判別式曲線上の線束 $\theta_N$ を
\[
 0\longrightarrow\mathcal O_{\mathbb P(N)}(-2)^3
 \xrightarrow{\lambda_0A_0+\lambda_1A_1+\lambda_2A_2}
 \mathcal O_{\mathbb P(N)}(-1)^3
 \longrightarrow i_*\theta_N\longrightarrow0
\]
で定めると、
\[
 \theta_N^{\otimes2}\simeq\mathcal O_{\Delta_N},\qquad
 H^0(\Delta_N,\theta_N)=0,\qquad
 (\pi_N)_*\mathcal O_{E_Q}\simeq\mathcal O_{\Delta_N}\oplus\theta_N.
\]
特に二つの種数一曲線のJacobianの間には自然な次数二の同種写像がある。
\end{theorem}

\begin{proof}
$\Delta_N$ が滑らかなので、その上で対称行列は階数二を持つ。
階数一以下なら随伴行列が零になり、行列式の全偏微分が消えるからである。
従って各退化二次曲線は相異なる二直線の和となる。

一つの直線が二つの独立なnetの二次曲線に含まれることはない。
そのような二次曲線があれば、それらのpencil全体が退化し、
$\Delta_N$ が射影直線を含むことになる。これは滑らかな三次曲線の既約性に反する。
よって $E_Q$ 上では $r_N$ の核が恒等的に一次元であり、その核を取ることにより
写像 $\pi_N$ が正則に定まる。その各ファイバーは退化二次曲線の二成分であり、
相異なる二点からなる。二次曲線の階数が二で一定であることから、局所的に
二変数非退化二次形式の等方直線を選ぶ被覆となり、有限エタール次数二である。

対称表示の双対性と $K_{\Delta_N}\simeq\mathcal O_{\Delta_N}$ から
$\theta_N\simeq\theta_N^*$ を得る。
表示のコホモロジーにより $H^0(\theta_N)=0$ なので、この二捩れは非自明である。
被覆の線束を同定するため、$\Delta_N$ 上で
\[
 \mathscr K=\ker\bigl(\mathcal O^3\longrightarrow\mathcal O(1)^3\bigr),
 \qquad \mathscr V=\mathcal O^3/\mathscr K
\]
と置く。$\mathscr V$ は $\mathcal O(1)$ に値を持つ階数二の非退化二次束である。
等方直線を選ぶ二重被覆のtrace零線束は、その判別式線束
\[
 \det(\mathscr V)\otimes\mathcal O(-1)
 =\mathscr K^*\otimes\mathcal O(-1)
\]
である。この記述は、局所的に二次形式を対角化して二つの等方方向を交換する作用を見れば分かる。
一方、対称性から上の行列の余核は $\mathscr K^*\otimes\mathcal O(1)$ なので、
表示を二つtwistすることにより
$\theta_N=\mathscr K^*\otimes\mathcal O(-1)$ となる。
これで被覆の線束の同定が得られる。
$H^0(\mathcal O_{\Delta_N}\oplus\theta_N)=\mathbb C$ により被覆は連結である。
\end{proof}

\subsection{古典的なnet写像との同定}
このモジュラー写像は\cite[Remark~4.1]{DolgachevRationality2008}のnet写像に一致する。
以下では関係ファイバーの補助曲線との対応を明示する。

\begin{corollary}[次数二同種写像のモジュライへの写像]
\label{cor:prym-y0two-v27}
一般の種数三Prym対から、非自明二捩れを持つ楕円曲線の粗モジュライへの支配的有理写像
\[
 \Psi:\mathcal R_3\dashrightarrow Y_0(2),\qquad
 (C,\eta)\longmapsto
 \bigl(\operatorname{Pic}^0(\Delta_N),\langle\theta_N\rangle\bigr)
\]
が得られる。補助曲線 $E_Q$ のJacobianは、
$\operatorname{Pic}^0(\Delta_N)$ を $\langle\theta_N\rangle$ で割った楕円曲線である。
\end{corollary}

\begin{proof}
対称最小分解の一意性により、同じ $(C,\eta)$ を表す $Q$ の選択差は
正ブロックの合同変換と平面座標変更である。
これらはnet $N$ とその対称一次表示を同型に移すので、
$(\Delta_N,\theta_N)$ は同型まで定まる。
Jacobianを用いることで、種数一曲線上の原点の選択を必要としない。
被覆の引戻し $\pi_N^*:\operatorname{Pic}^0(\Delta_N)\to\operatorname{Pic}^0(E_Q)$
の核は $\langle\theta_N\rangle$ であり、次数二の同種写像を与える。

滑らかさの開条件と非定数性は既存の一変数Gram族で確認できる。
その族では
\begin{align*}
 E_t&:\ 2tz_1^3+z_2^3+tz_3^3+(t-8)z_1z_2z_3=0,\\
 \Delta_t&:\ \lambda_0^3+4\lambda_1^3+t\lambda_2^3
                    -(t+4)\lambda_0\lambda_1\lambda_2=0.
\end{align*}
$t=1$ では両曲線が滑らかであり、$\det G=t-4\ne0$ である。
Hesse形の公式から
\[
 j(\Delta_t)=
 \frac{(t+4)^3\bigl((t+4)^3+864t\bigr)^3}
 {4t(t-2)^6(t+16)^3}
\]
を得る。これは非定数である。
十八次元のGram空間上の非定数な有理関数なので、四次Prym曲線が滑らかな
稠密開部分へ制限しても非定数性を保つ。
$Y_0(2)$ は既約な曲線であり、したがって $\Psi$ は支配的である。
\end{proof}

\paragraph{適用範囲.}
$E_Q$ の同型類はnetだけに依存するが、
線織面を定める階数二の核束 $\mathcal K_Q$ までnetだけに依存するとはここでは主張しない。
また、この次数二同種写像は補助種数一曲線のJacobianに関するものである。
種数三Prym対に付随する二次元Prym多様体の楕円因子を与えるという意味ではない。

\paragraph{再現入口.}
\nolinkurl{hermite_pfaffian_v27.py} は、二十一個の形式変数に関するPfaffianの小行列式表示と、
十八個の形式変数に関する相補小行列式の恒等式を検証する。
さらに独立な全座標例、行列式三の座標変更の共変性、上の二つのHesse曲線、
両者の $j$ 不変量が満たす次数二のモジュラー方程式を有理数上で検証する。
幾何学的被覆とその線束の同定は、有限計算に加えて上の証明を用いる。

\subsection{通常のPrym多様体との関係}
\label{subsec:classical-prym-context-v27}

同じnetから作る二つの被覆を整理しておく。
対称表示が定義される体上で、$C:ac-b^2=0$ に対応する非分岐二重被覆は
\[
 \widetilde C:\quad a(x)=r^2,\quad b(x)=rs,\quad c(x)=s^2
 \quad\subset\mathbb P^4
\]
という種数五の標準モデルを持つ。
この三つの二次超曲面のnetの判別式は
\[
 \Delta_N\ \cup\ \{4\lambda_0\lambda_2-\lambda_1^2=0\}
\]
という平面五次曲線である。
Bruinの構成\cite[Theorem 5.1, case 4]{BruinPrymGenusThree}では、一般の場合、
二次元Prym多様体は種数二曲線
\[
 F:\quad y^2=-\det(A_0+2tA_1+t^2A_2)
\]
のJacobianとして表される。その六つの分岐点は上の二次曲線と $\Delta_N$ の交点である。
一方、本稿の $E\to\Delta_N$ は、平面内の退化二次曲線の二直線成分を選ぶ
\emph{エタール}被覆である。
種数五の標準モデルの $W^1_4$ を使って判別式上に作る、六点で分岐する被覆とは異なる。
従って $\mathcal R_3\dashrightarrow Y_0(2)$ はnetに付随する補助不変量を記述している。

\section{等方Grassmann曲面の三成分境界}
\label{sec:grassmann-boundary-v28}

固定Gramファイバーを与える $X_Q$ には、有限二十一点配置の開部分から除かれる二つの平面がある。
ここでは閉包全体を決定する。基礎体は代数閉体、標数は零とし、$Q$ は一般とする。
前節と同じく $T=k^6$、$B=k^3$ と書き、
\[
 X=X_Q=\{W\in\operatorname{Gr}(4,T\oplus B):
       (\Omega_i\oplus\varepsilon_i)|_W=0\ (i=1,2,3)\}
\]
と置く。$\varepsilon_i$ は $\bigwedge^2B^*$ の標準基底である。
以下の三次曲線は $E=\{\operatorname{Pf}\Omega(z)=0\}$ である。

\begin{theorem}[三成分の単純正規交差曲面]
\label{thm:grassmann-snc-v28}
一般の $Q$ について、$X$ は被約な局所完全交叉曲面であり、
ちょうど三つの既約成分
\[
 X=R\cup\Pi_+\cup\Pi_-,\qquad \Pi_\pm\simeq\mathbb P^2
\]
を持つ。$R$ は $E$ 上の分解可能な射影直線束であり、ある
$\mathcal L\in\operatorname{Pic}^0(E)$ に対して
$R\simeq\mathbb P_E(\mathcal O_E\oplus\mathcal L)$ と書ける。
二つの平面は交わらず、
\[
 D_\pm=R\cap\Pi_\pm
\]
はそれぞれ平面内の滑らかな三次曲線である。
$D_\pm\to E$ は同型で、$R$ の互いに交わらない二つの切断をなす。
二重曲線の全ての点で、$X$ の完備局所環は $k[[a,b,c]]/(ab)$ である。
また
\[
 \omega_X\simeq\mathcal O_X,\qquad
 h^1(X,\mathcal O_X)=0,\qquad h^2(X,\mathcal O_X)=1.
\]
\end{theorem}

\subsection{射影の階数と二つの平面}

まず三つの $\Omega_i$ に共通等方的な三平面の零点スキームは、
一般には相補的な二つの被約点 $U_+,U_-$ である。
この開条件は $G=\operatorname{diag}(I_3,I_3)$ で確認できる。
このとき共通等方条件は
$v\times v'+w\times w'=0$ となる。
三平面の一方の射影が単射でなければ、他方の射影像は一次元以下となり、
その核の二つの独立なベクトルの外積が非零になることに矛盾する。
従って三平面は $w=Av$ というグラフであり、
$\bigwedge^2A=-I$、すなわち $A=\pm\sqrt{-1}I$ である。
外積二乗写像の微分は可逆なので二解は被約である。
相対Grassmann多様体の射影性とこのJacobi条件から、同じ性質を保つ開近傍が存在する。
これを $E$ の滑らかさと良い固定Gramファイバーの既存の非空開部分と交わらせる。

$T=U_+\oplus U_-$ に関して
\begin{equation}\label{eq:cross-block-M-v28}
 \Omega(z)=\begin{pmatrix}0&M(z)\\-M(z)^t&0\end{pmatrix},
 \qquad M(z)=z_1M_1+z_2M_2+z_3M_3
\end{equation}
と書ける。$E$ の方程式は符号を除いて $\det M(z)=0$ である。

$r=\operatorname{rank}(W\to B)$ とする。
$0\ne\ell\in W\cap B$ が存在すれば、等方条件から
$\ell\wedge\operatorname{pr}_B(W)=0$ なので $r\le1$ である。
従って $r\ge2$ なら $W\to T$ は単射である。
$r=2$ と仮定すると、その四次元像 $F\subset T$ への三形式の制限は、
二次元像 $\operatorname{pr}_B(W)$ の外積を通るため一次元以下を張る。
従ってパラメータ平面内のある直線に属する全ての $\Omega(z)$ が $F$ を零にする。
六次元の非退化交代形式に四次元等方部分空間は存在しないから、
この直線は $E$ に含まれる。これは $E$ の滑らかさに反する。
$r=0$ も同じ理由で起こらない。

$r=1$ なら $W\cap T$ は $U_+$ または $U_-$ である。
三形式に関する $U_\pm$ の共通直交補空間は $U_\pm$ 自身である。
実際、非退化な $\Omega(z)$ を一つ選べば、その形式に関して $U_\pm$ はLagrange部分空間である。
従って
\[
 W=U_\pm\oplus\ell,\quad \ell\in\mathbb P(B),
 \qquad
 \Pi_\pm=\{U_\pm\oplus\ell\}\simeq\mathbb P(B).
\]
残る $r=3$ の開部分を $X^\circ$ と書く。
ここでは $r\ge2$ だけから $W\cap B=0$ が従ったので、
前節の逆構成の証明を、有限点配置という追加条件を外した全ての $r=3$ の点に適用できる。
従ってこの開部分は、三つの誘導形式が基底となるスクロールの開部分の
エタール二重被覆そのものであり、滑らかである。

\subsection{境界における正確な局所方程式}

$U_+\oplus\ell_0$ の近傍で $B=\langle\ell_0\rangle\oplus B_2$ を選ぶ。
$W$ を、$u\in U_+$ に対応するベクトルと最後の一ベクトルによって
\[
 (u,Au,Cu),\qquad (0,d,\ell),
 \quad
 A:U_+\to U_-,\ C:U_+\to B_2,\ d\in U_-,\
 \ell\in\ell_0+B_2
\]
と表示する。$
\mathcal T_M(A)=(M_iA-A^tM_i^t)_{i=1}^3
$ は九次元空間の間の線形写像である。
これは共通等方三平面の零点スキームの $U_+$ における微分であり、
$U_+$ の被約性から可逆である。
最初の三ベクトル相互の等方条件は
\[
 \mathcal T_M(A)+(C^t\varepsilon_i C)_{i=1}^3=0
\]
なので、$A$ を $C$ の二次式として一意に消去できる。
最後のベクトルとの条件は
\[
 M_i d+C^t\varepsilon_i\ell=0\qquad(i=1,2,3)
\]
であり、$A$ を含まない。
写像 $B_2\to\bigwedge^2B$, $b\mapsto b\wedge\ell$ はその二次元像への同型なので、
この九本の式から $C$ を一意に消去できるための必要十分条件は
\begin{equation}\label{eq:boundary-local-kernel-v28}
 M(\ell)d=0
\end{equation}
である。従って $X$ はこのチャートで、二つの $\ell$ 座標と三つの $d$ 座標の中の
式\eqref{eq:boundary-local-kernel-v28}とスキームとして同型である。

$\det M(\ell)\ne0$ なら $d=0$ で、平面 $\Pi_+$ だけが残る。
$\det M(\ell)=0$ なら $M(\ell)$ は階数二である。
階数一以下なら余因子と三つの偏微分が全て消え、$E$ が特異になるからである。
二次単元ブロックを消去すれば、局所方程式は
\[
 f(\ell)\,v=0,\qquad f(\ell)=\det M(\ell)\text{ の単元倍}
\]
となる。$f=0$ は滑らかな三次曲線なので、二つの滑らかな分枝が横断的に交わる。
$\Pi_-$ でも同じ議論が成立する。
これにより $X$ は全ての点で純二次元の局所完全交叉であり、
平面以外の成分は滑らかである。

\subsection{主成分の線織構造と双対化層}

分解 $T=U_+\oplus U_-$ を用い、$\mathbb G_m$ を
\[
 g_\lambda=\operatorname{diag}(\lambda I_{U_+},\lambda^{-1}I_{U_-},I_B)
\]
によって $X$ に作用させる。これは三形式を保つ。
$g_{-1}$ のGrassmann多様体への作用は、残余符号 $\tau$ の作用に等しい。
スクロール上では、各核
$\ker\Omega(z)=\ker M(z)^t\oplus\ker M(z)$ の二直線を保ち、
それ以外の点に $\lambda^2$ による作用を与える。
二つの固有直線においては $F_k/k$ 上の三形式が一次独立でないので、
$X^\circ$ の像には含まれない。
一般の $z$ について、残る $\mathbb G_m$ 軌道は像に含まれる。
二重被覆の二点は $g_{-1}$ で交換されるため、その逆像は一つの自由な
$\mathbb G_m$ 軌道である。
従って $X^\circ$ は既約で、その $E$ 上の一般ファイバーは
分裂トーラスのtorsorである。体上の $\mathbb G_m$-torsorは自明なので、
その閉包 $R$ は $E\times\mathbb P^1$ と双有理である。
先の局所表示は、二つの平面に沿う主分枝がこの同じ $R$ に属することも示す。

滑らかな射影曲面 $R$ から $E$ への有理写像は射へ延長する。
実際、点のblow-upで不定点を解消すると、例外的な射影直線は種数一曲線へ定値に写るので降下できる。
この射を $p:R\to E$ と書く。
種数一曲線上の線織曲面の点のblow-upという表示から、垂直な既約曲線は全て有理曲線である。
従って種数一の二曲線 $D_\pm$ はともに $E$ を支配する。

Grassmann多様体の普遍部分束を $\mathcal U$ とする。
$X$ は階数十八の束 $(\bigwedge^2\mathcal U^*)^{\oplus3}$ の正則切断の零点であり、
\[
 \omega_{\operatorname{Gr}(4,9)}=\mathcal O(-9),\qquad
 \det((\bigwedge^2\mathcal U^*)^{\oplus3})=\mathcal O(9).
\]
随伴により $\omega_X\simeq\mathcal O_X$、従って
\begin{equation}\label{eq:anticanonical-boundary-v28}
 K_R+D_++D_-=0.
\end{equation}
一般の射影直線ファイバーと交差させると
$\deg(D_+\to E)+\deg(D_-\to E)=2$ である。
両次数は正なのでともに一となり、$D_\pm\to E$ は同型である。
これらを同一のパラメータ $z=\ell$ による同型とみなす必要はない。

局所表示から
\[
 N_{D_+/R}\simeq\ker(M|_E)(1),\qquad
 N_{D_-/R}\simeq\ker(M^t|_E)(1)
\]
を得る。ここでは境界平面の $\ell$ による三次曲線のパラメータ化を用いる。
核は $\mathcal O_E^3\to\mathcal O_E(1)^3$ の核である。
一次行列式表示の余核線束 $\mathcal A$ は
$0\to\mathcal O(-1)^3\to\mathcal O^3\to i_*\mathcal A\to0$
から次数三を持つ。これを $E$ へ制限して得るTor列により、
$\ker(M|_E)\simeq\mathcal A(-2)$ は次数 $-3$ である。
転置でも同様なので、二つの法束の次数は零となる。

$R$ から相対極小線織曲面へ向けて $(-1)$ 曲線を収縮する。
式\eqref{eq:anticanonical-boundary-v28}と随伴から、各 $(-1)$ 曲線は
二つの切断の和と交差数一を持つ。
従って一方の切断だけと横断的に交わり、収縮は二つの切断の非交差性を保つ。
一回ごとに自己交差数の和は一だけ増える。
初めの和は零であり、極小線織面の互いに交わらない二つの切断の自己交差数の和も零であるから、
収縮する曲線は存在しない。
従って $R$ 自身が射影直線束であり、二つの切断によって分解可能で、対応する線束の次数は零である。

最後に、三成分の構造層の完全列
\[
 0\to\mathcal O_X\to
 \mathcal O_R\oplus\mathcal O_{\Pi_+}\oplus\mathcal O_{\Pi_-}
 \to\mathcal O_{D_+}\oplus\mathcal O_{D_-}\to0
\]
を用いる。$H^1(R,\mathcal O_R)\to
H^1(D_+,\mathcal O_{D_+})\oplus H^1(D_-,\mathcal O_{D_-})$
は一次元から二次元への単射である。
従って $h^1(\mathcal O_X)=0$、$h^2(\mathcal O_X)=1$ を得る。
これで定理を証明した。

\subsection{一般平滑化の全空間に現れる十八個の普通二重点}

\begin{corollary}[平滑化と二重曲線の九点ずつの零点]
\label{cor:eighteen-smoothing-nodes-v28}
三つの交代形式を $T\oplus B$ 上の一般の三形式へ一次に摂動すると、
$X$ を特殊ファイバー、種数二十の滑らかな偏極K3曲面を一般ファイバーとする平坦族を得る。
この一般一次摂動の全空間は、特殊ファイバー上にちょうど十八個の三次元普通二重点を持つ。
それらは二つの二重曲線に九個ずつ存在する。
\end{corollary}

\begin{proof}
正則切断の零点を摂動した族は、$0$ の近傍で局所完全交叉かつ平坦である。
一般の三つの交代形式の零点が種数二十のK3曲面となることには
Mukaiの構成\cite[Theorem~5.1]{MukaiGenus20}を用いる。
一般の摂動方向はその滑らかな開部分と交わる。

単純正規交差曲面の局所方程式 $ab=0$ から、局所一次変形の層は
\[
 \mathcal T_X^1|_{D_\pm}
   \simeq N_{D_\pm/R}\otimes N_{D_\pm/\Pi_\pm}
   =N_{D_\pm/R}\otimes\mathcal O_{D_\pm}(3)
\]
であり、その次数は九である。
全ての交代形式の摂動は、束
$(\bigwedge^2\mathcal U^*)^{\oplus3}$ の大域切断を与える。
この束は大域生成され、切断の変形から $\mathcal T_X^1$ への写像は各点で全射なので、
得られる一次平滑化切断の線形系は二重曲線の上で基点を持たない。
従って一般の切断は、それぞれ九個の単純零点を持つ。

二重曲線上の局所座標を $c$、基底パラメータを $t$ とすると、全空間の方程式は
\[
 ab-t h(c)+O(ta,tb,t^2)=0
\]
と書ける。$h(c)\ne0$ では全空間は滑らかである。
$h$ の単純零点では二次項が非退化な四変数二次形式となるので、三次元普通二重点である。
特殊ファイバーの滑らかな点では全空間も滑らかであり、従って、これが特殊ファイバーに沿う全ての特異点である。
他の特異ファイバーを除くよう基底を $0$ の周囲で縮めれば、全空間の特異点もこの十八点だけとなる。
\end{proof}

この族の全空間は十八点で特異である。
従って、追加の双有理修正を行わずに滑らかな全空間を持つ半安定モデルとは呼ばない。
二次元の $X$ の正規交差特異点と、三次元の全空間の普通二重点は区別する。
ここで数えた十八点は二重曲線上の平滑化切断の零点であり、六次Prym族の $D_{\rm mov}$ とは別の構成である。

\subsection{二平面の小解消と射影的Kulikovモデル}
\label{subsec:grassmann-kulikov-v28}

前項の十八個の普通二重点は、二つの平面を使って同時に解消できる。
ここでいうKulikovモデルは、滑らかな全空間、被約単純正規交差特殊ファイバー、
自明な相対標準束を持つK3曲面の一径数退化を意味する
\cite{KulikovDegenerations,PerssonPinkhamDegenerations}。
以下では基底の $0$ を含む開近傍に制限し、他の特異ファイバーがあれば除いておく。

\begin{theorem}[平面を用いた射影的小解消]
\label{thm:grassmann-kulikov-v28}
系\ref{cor:eighteen-smoothing-nodes-v28}の一般一次摂動族を
$\mathscr X\to S$ とする。
特殊ファイバーの二成分 $\Pi_+,\Pi_-$ を $\mathscr X$ のWeil因子として扱い、
そのイデアルを順次blow-upすると、射影的な小解消
\[
 \widetilde{\mathscr X}\longrightarrow\mathscr X
\]
を得る。これはcrepantであり、$\widetilde{\mathscr X}$ は滑らかで
$\omega_{\widetilde{\mathscr X}/S}\simeq\mathcal O$ である。
特殊ファイバーは
\[
 \widetilde X=\widetilde\Pi_+\cup R\cup\widetilde\Pi_-,\qquad
 \widetilde\Pi_\pm\simeq
 \operatorname{Bl}_{Z_\pm}\mathbb P^2
\]
という三成分の被約単純正規交差曲面である。
ここで $Z_\pm$ は前項の平滑化切断の九個の相異なる零点であり、
$R$ および二重曲線 $D_\pm$ は同型を保つ。
この族は、両端が有理曲面、中間が楕円線織面であるType~IIのKulikovモデルとなる。
基底変換は必要ない。
\end{theorem}

\begin{proof}
普通二重点の近傍では、特殊ファイバーの二成分を $a=0$ と $b=0$ として、
全空間の方程式は
\[
 ab-tH(a,b,c,t)=0,\qquad H_c(0)\ne0
\]
と書ける。$c$ を $H$ で置き換えれば、基底パラメータ $t$ と二成分を保ったまま
\begin{equation}\label{eq:node-small-resolution-v28}
 ab=tc
\end{equation}
という局所形になる。
平面側を $\{a=t=0\}$ とし、そのイデアル $(a,t)$ をblow-upする。
$a$ が生成元となるチャートでは
\[
 t=au,\quad b=uc
\]
であり、$(a,u,c)$ が滑らかな三次元座標を与える。
特殊ファイバーは $au=0$ である。
$t$ が生成元となるチャートでは
\[
 a=tv,\quad c=vb
\]
となり、$(t,v,b)$ が滑らかな座標を与える。
この二チャートを合わせると、元の普通二重点だけが一本の射影直線に置き換わる。
従ってこれは小解消であり、解消後の特殊ファイバーは被約かつ単純正規交差である。

選択した成分 $\{a=t=0\}$ への制限は、その局所座標 $(b,c)$ の原点のblow-upである。
実際、上の二チャートでは $b=uc$ と $c=vb$ がその二つの標準チャートとなる。
他方の成分 $\{b=t=0\}$ は $(a,c)$ を座標とする $u=0$ のチャートで同型に保たれる。
従って平面を選んで解消すれば、平面だけが九点でblow-upされ、$R$ は変わらない。

全空間の普通二重点以外では、各 $\Pi_\pm$ はCartier因子である。
平面の内部では特殊ファイバー自身を局所的に定義する因子であり、
二重曲線の平滑化切断が非零の点では $ab=t\cdot\text{unit}$ より一つの式で定義される。
従ってこのblow-upは普通二重点の外で同型である。
二平面は交わらないため、それぞれのイデアルのblow-upは互いの中心に影響しない。
これにより、局所小解消は大域的かつ射影的な二回のblow-upとして実現する。

$\mathscr X$ はGrassmann多様体と基底の積の中で、
$\bigl(\bigwedge^2\mathcal U^*\bigr)^{\oplus3}$ の正則切断により定義される。
相対随伴により $\omega_{\mathscr X/S}\simeq\mathcal O$ である。
小解消には例外因子がなく、$\mathscr X$ の標準束は可逆なので、
標準因子の差は零となる。従って解消はcrepantであり、
$\omega_{\widetilde{\mathscr X}/S}\simeq\mathcal O$ を得る。
一般ファイバーは変更されず、種数二十のK3曲面のままである。
上のチャートが示す $au=t$ 型の局所形により、解消後の族は半安定である。
特殊ファイバーは有理曲面、楕円線織面、有理曲面の鎖となり、Type~IIの形を持つ。
\end{proof}

\begin{corollary}[二重曲線の法束の相殺]
\label{cor:normal-cancellation-v28}
解消後の二重曲線について
\[
 N_{D_\pm/\widetilde\Pi_\pm}
 \simeq N_{D_\pm/R}^{-1},\qquad
 \mathcal T^1_{\widetilde X}|_{D_\pm}\simeq\mathcal O_{D_\pm}
\]
が成立する。両端の有理曲面では $D_\pm$ は自己交差数零の反標準因子である。
\end{corollary}

\begin{proof}
摂動の一次平滑化切断を $h_\pm$ とする。
その零除数は $Z_\pm$ であり、前項から
\[
 \mathcal O_{D_\pm}(Z_\pm)
 \simeq N_{D_\pm/R}\otimes\mathcal O_{D_\pm}(3)
\]
である。滑らかな三次曲線上の九点をblow-upすると、その法束は
\[
 N_{D_\pm/\widetilde\Pi_\pm}
 =\mathcal O_{D_\pm}(3)(-Z_\pm)
 =N_{D_\pm/R}^{-1}
\]
となる。これが局所一次変形層の自明性も与える。
また $\operatorname{Bl}_9\mathbb P^2$ 上では、三次曲線の狭義変換は
$3H-\sum_{j=1}^9E_j=-K$ であり、その自己交差数は $9-9=0$ である。
\end{proof}

\paragraph{偏極と有理端成分の意味.}
Pl\"ucker偏極の引戻しは、解消後も半豊富で、一般ファイバー上で次数三十八の偏極を与える。
例外的な十八本の射影直線上では次数零となる。
また、端成分が九点blow-upで反標準曲線の自己交差数が零であることだけから、
反標準線形系が楕円pencilをなすとはいえない。
ここでは二つの端を有理曲面と記述し、法束の自明性や捩れ性を追加で仮定しない。

\section{行列式線束、境界の平行移動と端成分の半豊富性}
\label{sec:ruled-line-bundle-v29}

この節では、前節の主成分を定める次数零の線束を同定する。
基礎体は標数零の代数閉体とし、三つの交代形式がちょうど二つの被約で相補的な
共通等方三平面 $U_+,U_-$ を持ち、Pfaffian三次曲線が滑らかで、
定理\ref{thm:grassmann-snc-v28}の三成分境界が成立する場合を扱う。
二平面を順序付け、$T=U_+\oplus U_-$ に関して
\[
 \Omega(z)=\begin{pmatrix}0&M(z)\\-M(z)^t&0\end{pmatrix},
 \qquad E=\{\det M=0\}\subset\mathbb P^2,
 \qquad H=\mathcal O_E(1)
\]
と書く。$M$ は $3\times3$ の一次行列である。
行列式線束 $\theta$ を
\begin{equation}\label{eq:determinantal-theta-v29}
 0\longrightarrow\mathcal O_{\mathbb P^2}(-2)^3
 \xrightarrow{M}\mathcal O_{\mathbb P^2}(-1)^3
 \longrightarrow i_*\theta\longrightarrow0
\end{equation}
により定める。この線束は、netの判別曲線 $\Delta_N$ 上の二位数線束 $\theta_N$ と区別する。
$\theta_N$ がHermite被覆 $E\to\Delta_N$ を与えるのに対し、
$\theta$ は $E$ 上の次数零の線束であり、一般には有限位数でない。

\begin{theorem}[主成分と符号商を定める線束]
\label{thm:ruled-determinantal-line-v29}
上の仮定の下で、$\theta\in\operatorname{Pic}^0(E)$ は非自明であり、
\begin{equation}\label{eq:kernel-theta-v29}
 \mathcal K=H^{-1}\otimes(\theta^{-1}\oplus\theta).
\end{equation}
直線をパラメータ化する射影化の規約では、主成分と符号商スクロールは
\begin{equation}\label{eq:ruled-square-v29}
 R\simeq\mathbb P_E(\mathcal O_E\oplus\theta),\qquad
 R/\langle\tau\rangle\simeq
 S=\mathbb P_E(\mathcal K)\simeq
 \mathbb P_E(\mathcal O_E\oplus\theta^2)
\end{equation}
となる。商写像は各ファイバーで二乗写像であり、二つの核固有切断に沿って分岐する。
また
\begin{equation}\label{eq:normal-theta-v29}
 N_{D_+/R}\simeq\theta,\qquad N_{D_-/R}\simeq\theta^{-1}.
\end{equation}
境界平面の三次曲線を元のパラメータ $x=\ell$ で $E$ と同定すると、
線織射影 $p:R\to E$ の二切断への制限は
\begin{equation}\label{eq:boundary-translation-v29}
 p|_{D_+}=t_\theta,\qquad p|_{D_-}=t_{\theta^{-1}}
\end{equation}
である。ここで $t_\theta$ は
$\theta\simeq\mathcal O_E(t_\theta(x)-x)$ によって定まる
$\operatorname{Jac}(E)$ の平行移動であり、$E$ の原点を選ばずに定義される。
さらに、二共通平面の被約性の仮定から $\theta^2\not\simeq\mathcal O_E$ である。
\end{theorem}

\subsection{核束と境界法線}

$E$ が滑らかなので、$M$ は $E$ 上の全ての点で階数二を持つ。
従って\eqref{eq:determinantal-theta-v29}の余核は $E$ 上の線束である。
Hilbert多項式からその次数は零であり、負の次数の射影平面のコホモロジーから
$H^0(E,\theta)=0$ を得る。従って $\theta$ は非自明である。
これは一次行列式表示の古典的対応\cite[Proposition~3.1]{BeauvilleDeterminantal}の
三次曲線の場合に当たる。
同じ分解を $\mathcal O_{\mathbb P^2}(-3)$ で双対化すると
\[
 \mathcal Ext^1_{\mathbb P^2}(i_*\theta,\mathcal O(-3))
 \simeq i_*(\theta^{-1})
\]
なので、転置行列の線束は $\theta^{-1}$ となる。

式\eqref{eq:determinantal-theta-v29}を一回捻って $E$ へ制限する。
Cartier因子 $E$ の次数が三であることから、そのTor項は $\theta H^{-2}$ である。
さらに一回捻ると
\[
 \ker(M:\mathcal O_E^3\to H^{\oplus3})=\theta H^{-1},\qquad
 \ker(M^t:\mathcal O_E^3\to H^{\oplus3})=\theta^{-1}H^{-1}.
\]
これが\eqref{eq:kernel-theta-v29}を与える。
境界チャートの式 $M(\ell)d=0$ による法線表示から
\eqref{eq:normal-theta-v29}も従う。
ただし、この段階の法線表示では境界側のパラメータ $\ell$ を用いている。
次項でその同型と線織基底の同型の差を計算する。

\subsection{二つのパラメータの差}

\begin{lemma}[行列式表示の平行移動]
\label{lem:determinantal-translation-v29}
$x\in E$ と $0\ne d\in\ker M(x)$ を取る。
$u^tM_i d=0$ $(i=1,2,3)$ を満たす $u$ はスカラーを除いて一意に定まり、
唯一の $y\in E$ に対して $u\in\ker M(y)^t$ となる。そのとき
\[
 \theta\simeq\mathcal O_E(y-x).
\]
\end{lemma}

\begin{proof}
線形写像 $z\mapsto M(z)d$ は階数二である。
階数一以下なら、その共通零点集合が平面内の直線を含み、
$\det M=0$ もその直線を含んでしまう。
従って三つの一次形式 $M(z)d$ の共通零点は $x$ ただ一つである。
二つの独立な一次形式の $E$ への制限の共通零因子は単純な $x$ となる。

右核と左核による写像は、それぞれ完全線形系
$|H\theta^{-1}|$ と $|H\theta|$ による平面三次曲線の埋め込みである。
実際、余核の分解を一回捻った完全列とその双対は、
三つの座標切断がこれらの次数三の線束の大域切断全体を張ることを示す。
genus one上の次数三の線束は非常に豊富である。

$u^tM(z)$ は固定の右関係 $d$ を持つので、その $z$ に関する係数行列は特異である。
この行列式曲線は左核写像の像に等しい。
行列式が恒等的に零なら全ての $u\in\mathbb P(U_+)$ が左核となるが、
対応するincidenceは $E$ と同型な一次元空間であるから不可能である。
従って $u$ は唯一の $y\in E$ の左核となる。

$\mathcal I=\operatorname{im}(M:\mathcal O_E^3\to H^{\oplus3})$ と置くと
$\det\mathcal I=H\theta^{-1}$ である。
固定 $d$ による切断を飽和させた部分線束は $\mathcal O_E(x)$ であり、
固定 $u$ とのペアリングによる $\mathcal I\to H$ の像は $H(-y)$ である。
両者の合成は零であり、部分線束が飽和しているので
\[
 0\longrightarrow\mathcal O_E(x)\longrightarrow\mathcal I
 \longrightarrow H(-y)\longrightarrow0
\]
を得る。行列式を取れば
$H\theta^{-1}\simeq\mathcal O_E(x)\otimes H(-y)$ であり、結論が従う。
\end{proof}

境界チャートでは $d$ が法線方向、$C$ がその一次式、$A$ が二次式であった。
境界から離れた点の $W\cap T$ は $\ker C$ のグラフであり、
補題の $u$ が $\ker C$ を張る。
従って境界での極限核直線は $u\in\ker M(y)^t$ であって、
そのスクロールパラメータは $y=t_\theta(x)$ となる。
転置行列を用いればもう一方は $t_{\theta^{-1}}$ である。
これで\eqref{eq:boundary-translation-v29}を得る。
平行移動は $\operatorname{Pic}^0(E)$ に恒等的に作用するので、
法線束の式は元の線織基底でも同じ $\theta^{\pm1}$ となる。
互いに交わらない二切断を持つ射影直線束は分解可能であり、
一方の切断の法線が $\theta$ であることから $R\simeq\mathbb P_E(\mathcal O_E\oplus\theta)$ を得る。

\subsection{残余符号商の完備化}

階数の分類により、$\operatorname{rank}(W\to B)=3$ の開部分は
$X^\circ=R\setminus(D_+\cup D_-)$ に等しい。
スクロールの逆構成は、その符号商を $S$ の核固有二切断を除いた開部分に同型化する。
ここでは有限点配置という追加の開条件を課していない。
$R$ の各ファイバーから二切断を除いた部分は非空な一つの $\mathbb G_m$ 軌道である。
$S$ の各ファイバーから核固有二点を除いた部分も一つの軌道であり、
写像はトーラス作用と可換である。
従って商の像は、全てのファイバーでこの軌道全体に等しい。

$\tau$ はトーラスの $-1$ に当たり、各ファイバーでの商は二乗写像となる。
遷移関数も二乗されるため、この写像は二切断を加えた射影直線束へ
\[
 \mathbb P_E(\mathcal O_E\oplus\theta)
 \longrightarrow\mathbb P_E(\mathcal O_E\oplus\theta^2)
\]
と延長する。これが\eqref{eq:ruled-square-v29}を証明する。

最後に $\theta^2\simeq\mathcal O_E$ と仮定すると、対称行列式表示の理論
\cite[Theorem~B and Corollary~2.4]{BeauvilleDeterminantal}により
$M$ は左右の基底変換で対称行列にできる。
その場合、全ての $a\in k$ に対して $\{(u,au):u\in k^3\}$ が共通等方三平面となる。
これは二つの被約点という仮定に反する。定理の証明が完了した。

\subsection{有限位数とHalphen指数}

\begin{proposition}[有理端成分の半豊富性の判定]
\label{prop:halphen-torsion-v29}
定理\ref{thm:grassmann-kulikov-v28}の端成分を
$Y=\operatorname{Bl}_9\mathbb P^2$、二重曲線を $D$ とする。
このとき $D=-K_Y$、$D^2=0$ であり、$\mathcal O_D(D)$ は
$\theta$ または $\theta^{-1}$ に同型である。
$-K_Y$ が半豊富であるための必要十分条件は $\theta$ が有限位数であることである。
位数が $m$ なら $|mD|$ が最初の基点を持たないgenus-one pencilであり、
$Y$ はHalphen指数 $m$ を持つ。
位数が無限なら
\[
 h^0(Y,\mathcal O_Y(nD))=1\quad(n\ge1),
 \qquad\kappa(-K_Y)=0,\quad\nu(-K_Y)=1.
\]
この場合 $-K_Y$ はnefだが半豊富でない。
\end{proposition}

\begin{proof}
反標準因子と法線の同定は、系\ref{cor:normal-cancellation-v28}と
\eqref{eq:normal-theta-v29}から従う。
$D$ 以外の既約曲線とは非負に交わり、$D^2=0$ なので $D$ はnefである。
\[
 0\longrightarrow\mathcal O_Y((n-1)D)\longrightarrow
 \mathcal O_Y(nD)\longrightarrow\mathcal O_D(nD)\longrightarrow0
\]
を用いる。$\theta$ が無限位数なら最後の線束は常に非自明な次数零の線束なので
大域切断を持たず、帰納的に $h^0(nD)=1$ を得る。

位数が $m$ なら $0\le j<m$ で $h^0(jD)=1$ である。
Riemann--Rochから $\chi(\mathcal O_Y(jD))=1$、
また $H^2=H^0(-(j+1)D)^*=0$ なので $H^1=0$ となる。
$n=m$ の完全列は $h^0(mD)=2$ を与え、その新しい切断は $D$ 上で非零定数になる。
もう一方の切断の零因子は $mD$ なので、このpencilに基点はない。

射 $Y\to\mathbb P^1$ のStein分解の中間曲線は、$Y$ が有理曲面なので $\mathbb P^1$ である。
有限部分の次数が $e>1$ なら、一般の連結ファイバー $F$ に対して $mD=eF$ となる。
$\operatorname{Pic}(Y)$ は自由で $D$ は原始的だから $e\mid m$、
$F=(m/e)D$ であり、$m$ の最小性に反する。
従って一般ファイバーは連結であり、標数零での一般滑らかさと随伴から種数一である。
\end{proof}

\begin{corollary}[幾何学的一般点の端成分]
\label{cor:generic-nonsemiample-ends-v29}
逆Gram族の幾何学的一般点、および複素数体上の非常に一般の点では、
$\theta$ は無限位数である。
従って上のType~II Kulikovモデルの二つの有理端成分では反標準束は半豊富でない。
また $R$ と $S$ の線織面としての同型類はnet $N$ のみからは決まらない。
\end{corollary}

\begin{proof}
別稿\cite{PaperVI}の逆Gram族の支配性により、
$(E,\{\theta,\theta^{-1}\})$ は二次元の幾何学的moduliを支配する。
各 $n\ge1$ について $\theta^n=\mathcal O_E$ は真の閉条件であり、
幾何学的一般点はそのいずれにも属さない。
複素数体上では、これら可算個の真閉部分の合併の外で同じ結論が成立する。
netのmoduliは一次元で $E$ はnetのみから決まるから、
二次元を支配する $\theta$ はnetだけの不変量ではない。
二乗写像は有限なので、$S$ に現れる $\theta^2$ についても同様である。
\end{proof}

この系は、ある一つの非空Zariski開集合の全ての複素点で無限位数になるとは述べていない。
有限位数の場合の判定は命題\ref{prop:halphen-torsion-v29}による。
二つの被約共通平面を持つ開部分では $\theta^2\ne\mathcal O_E$ なので、
有限位数が現れる場合のHalphen指数は少なくとも三である。
全ての指数がこの逆Gram族の良い開部分で実現されるかは、ここでは主張しない。

\subsection{整数格子、偏極、および $A_8\oplus A_8$ の根系}
\label{subsec:kulikov-lattice-v29}

本項では複素数体上で、定理\ref{thm:grassmann-kulikov-v28}のKulikovモデルを考える。
両端を $Y_+,Y_-$、中間成分を $R$ と書く。
$Y_\pm$ 上の直線の引戻しを $H_\pm$、九本の例外曲線を $E_i^\pm$ とし、
\[
 d_\pm=3H_\pm-\sum_{i=1}^9 E_i^\pm
\]
と置く。$R$ 上では一方の二重曲線の数値類を $S$、線織のファイバー類を $F$ とする。
二つの二重曲線の法束は次数零なので、両者は同じ数値類 $S$ を持ち、
$S^2=F^2=0$、$SF=1$ である。

正規化の二次コホモロジー格子は
\[
 \Lambda=I_{1,9}\oplus U\oplus I_{1,9},\qquad
 I_{1,9}=\langle H,E_1,\ldots,E_9\rangle,
 \quad H^2=1,\ E_iE_j=-\delta_{ij}
\]
となる。ここで $U$ は基底 $S,F$ の双曲平面である。
次の二つのベクトルを用いる。
\[
 v_+=(d_+,-S,0),\qquad v_-=(0,S,-d_-),\qquad
 J=\mathbb Z v_+\oplus\mathbb Z v_-.
\]
二重曲線上での次数が一致する条件は $J^\perp$ であり、
成分を垂直因子として動かして得る類は $J$ を張る。
この数値Cartier格子の記述はType II退化の標準的な格子記述の三成分版である
\cite[\S2]{KondoTypeII}。

\begin{theorem}[偏極付き極限格子]
\label{thm:kulikov-lattice-v29}
上のKulikov退化について、次が成立する。
\begin{enumerate}
\item 重み二の整数格子は
\[
 L:=J^\perp/J\simeq U\oplus E_8(-1)\oplus E_8(-1)
\]
である。これは近傍のK3格子における、モノドロミー像の飽和 $I$ に対する
$I^\perp/I$ と同型である。
\item Pl\"ucker束の極限類は
\[
 h=(H_+,6S+3F,H_-)\pmod J,
 \qquad h^2=38
\]
であり、原始的である。その直交格子 $L_h=h^\perp\subset L$ は
階数十七の負定値偶格子で、行列式の絶対値は三十八、
判別群と判別形式は
\[
 A_{L_h}\simeq\mathbb Z/38,\qquad q_{L_h}(1)=-\frac1{38}\pmod{2\mathbb Z}
\]
である。
\item $L_h$ の平方 $-2$ のベクトルは、ちょうど
\[
 E_i^+-E_j^+\quad\text{および}\quad E_i^--E_j^-
 \qquad(i\ne j)
\]
である。従って根系は $A_8\oplus A_8$、根の総数は百四十四である。
\item モノドロミーの対数 $N$ は $N^2=0$、階数二であり、
その最大整数因子は二である。向きと $I$ の原始基底 $\delta,\mu$ を選ぶと
\[
 N(x)=2\bigl((x\cdot\mu)\delta-(x\cdot\delta)\mu\bigr),
 \qquad \operatorname{Im}N=2I
\]
と書ける。従って群としての指数は $[I:\operatorname{Im}N]=4$ である。
\end{enumerate}
\end{theorem}

\begin{proof}
$J$ は等方的である。また $E_1^+,E_1^-$ の係数で取る二次小行列式が $\pm1$ なので原始的である。
整数上の同定も確認する。Mayer--Vietoris列で
$H^1(R,\mathbb Z)\to H^1(D_+,\mathbb Z)\oplus H^1(D_-,\mathbb Z)$ は
二つの同型の差である。その余核と次数の整合条件から
\[
 0\longrightarrow\mathbb Z^2\longrightarrow H^2(\widetilde X,\mathbb Z)
 \longrightarrow J^\perp\longrightarrow0
\]
を得る。Clemens収縮による近傍のK3ファイバーへの特殊化は、左の項を重み一の飽和 $I$ へ、
全体を $I^\perp$ へ写す。垂直因子の類 $J$ は消えるので、積を保つ整数準同型
$L\to I^\perp/I$ が得られる。Clemens--Schmid列によりこれは有理数上で同型である
\cite{FriedmanTorelli,KondoTypeII}。
以下の明示基底は $L$ がユニモジュラーであることを示す。
一方、K3格子はユニモジュラーで $I$ は原始等方二平面なので $I^\perp/I$ もユニモジュラーである。
従ってこの準同型の像の有限指数は一であり、整数上でも同型となる。

さらに格子の同型を明示する。
\[
 e=(E_9^+,S+F,E_9^-),\qquad f=(d_+,0,0)
\]
は $J^\perp$ に属し、商で $e^2=f^2=0$、$ef=1$ である。
各端で
\[
 \alpha_i^\pm=E_i^\pm-E_{i+1}^\pm\ (1\le i\le7),\qquad
 \alpha_8^\pm=H_\pm-E_1^\pm-E_2^\pm-E_3^\pm
\]
を取る。これらは $e,f$ および他方の端の根と直交し、各組のGram行列は負の $E_8$ 行列である。
得られる階数十八の部分格子の行列式は $-1$ なので、$J^\perp/J$ 全体を張る。
これにより偶性とユニモジュラー性も直接分かる。

Pl\"ucker束は両端で $H_\pm$、二重曲線上で次数三である。
$R$ 上で $aS+bF$ と書けば $b=3$ であり、自己交差三十六から $a=6$ となる。
従って上の $h$ の式を得る。
その平方は $1+36+1=38$ である。
$h\alpha_8^+=1$ なので $h$ は原始的で、偶ユニモジュラー格子内の
原始ベクトルの直交補の判別形式は $\langle38\rangle$ の判別形式の反対となる。
Hodge指数型の署名計算から $L_h$ は負定値である。

平方 $-2$ の類 $x\in L_h$ の代表を
\[
 x=\left(aH_+-\sum b_iE_i^+,\ cS+dF,
             a'H_--\sum b'_iE_i^-\right)
\]
と書く。$x\in J^\perp$ と $xh=0$ は
\[
 \sum b_i=3a-d,\qquad \sum b_i'=3a'-d,\qquad
 a+a'+3c+6d=0
\]
と同値である。Cauchyの不等式を各九個の係数に用いると
\begin{align*}
 -x^2
 &=\sum b_i^2+\sum (b'_i)^2-a^2-(a')^2-2cd\\
 &\ge \frac{(3a-d)^2+(3a'-d)^2}{9}-a^2-(a')^2
            +\frac{2d(a+a'+6d)}3
 =\frac{38d^2}{9}.
\end{align*}
$-x^2=2$ なので $d=0$ である。
$v_+,v_-$ の整数倍を加え、$a,a'\in\{0,1,2\}$ としてよい。
このとき $a+a'\equiv0\pmod3$ である。
$a=1$ または $2$ なら、$\sum b_i=3a$ と整数性から
$\sum b_i^2-a^2\ge2$ となる。
従って $(a,a')=(1,2),(2,1)$ は $-x^2\ge4$ を与え、根ではない。
残る場合は $a=a'=c=d=0$ であり、係数の和は各端で零、平方和の合計は二である。
従って一方の端で二係数が $1,-1$、他が全て零である。
これが記載した全ての根である。

最後に、Type IIのモノドロミー対数が
$N(x)=(x\lambda)\delta-(x\delta)\lambda$ と書け、その非原始性指数が
二重曲線の本数に等しいという標準結果を使う
\cite[Theorem 3.11]{AlexeevEngelHanAutomorphisms}。
本例の二重曲線は二本なので、$\lambda=2\mu$ と取れる。
$I$ はK3格子の原始等方二平面であり、K3格子はユニモジュラーなので、
それとの対の写像は $I^*$ へ全射である。従って $\operatorname{Im}N=2I$ を得る。
\end{proof}

\paragraph{根格子と全境界格子の区別.}
根が張る部分格子 $A=A_8(-1)^{\oplus2}\subset L_h$ の行列式は八十一である。
その飽和 $\overline A$ は $[\overline A:A]=3$、行列式九を持つ。
実際、根の有理線形包は上の表示で $d=0$ の部分であり、
$a,a'$ の剰余は $(0,0),(1,2),(2,1)$ の三つとなる。
その直交補の原始生成元は
\[
 v=(3H_+,-20S+9F,3H_-)\pmod J,\qquad v^2=-342
\]
である。よって
\[
 [L_h:A\oplus\mathbb Z v]=27,\qquad
 [L_h:\overline A\oplus\mathbb Z v]=9.
\]
従って根系 $A_8\oplus A_8$ だけで全格子を記述したことにはならない。
定理の構成と付属の整数Gram行列が、偏極付き境界格子そのものを指定する。
これは極限の格子であり、一般K3ファイバーのPicard格子が階数十七になるという主張ではない。

\paragraph{再現可能な有限部分.}
\nolinkurl{kulikov_lattice_v29.py} は、二つの等方生成元、十八本の商基底、
Pl\"ucker類、直交補の整数Gram行列、根格子の飽和指数を有理整数上で検算する。
さらに独立な有理数LDL列挙により、平方 $-2$ の全百四十四根を確認する。
Type IIのモノドロミー公式とコホモロジーの同定は上の幾何学的議論による。

\section{十八個の零点から得られる偏極付き境界period}
\label{sec:period-realization-v30}

本節では複素数体上で議論する。前節までの一般の逆Gramパラメータから得られる
Type~IIモデル
\[
 \widetilde X=Y_+\cup R\cup Y_-,\qquad
 Y_\pm=\operatorname{Bl}_{Z_\pm}\mathbb P^2,
 \qquad R=\mathbb P_E(\mathcal O_E\oplus\theta)
\]
を考える。$E$ は滑らかな平面三次曲線、$H=\mathcal O_E(1)$、
$\theta\in\operatorname{Pic}^0(E)$ は
$\theta^{\otimes2}\not\simeq\mathcal O_E$ を満たす。
ここでいう一般性は、既に構成した三成分の単純正規交差モデルが成立する
空でないZariski開集合を指す。

本節の主張は、二つの九点配置が動くというだけでなく、
その変化が偏極付きType~II periodの全方向を与えることである。
正確なperiodの受け皿は
\[
 \operatorname{Hom}(L/\mathbb Zh,\operatorname{Pic}^0(E))
\]
である。$h^\perp\subset L$ への制限とは有限指数の違いがあるため、
両者を区別する。

\subsection{一次摂動は二つの完全線形系を独立に動かす}

$T=U_+\oplus U_-$ に関して
\[
 \Omega(z)=\begin{pmatrix}0&M(z)\\-M(z)^t&0\end{pmatrix},
 \qquad
 0\longrightarrow\mathcal O_{\mathbb P^2}(-2)^3
 \xrightarrow{M}\mathcal O_{\mathbb P^2}(-1)^3
 \longrightarrow i_*\theta\longrightarrow0
\]
と書く。行列式表示と次数零の非自明線束との対応自体は古典的であり、
ここでは \cite[Propositions~3.1, 3.3]{BeauvilleDeterminantal} の形を用いる。
$D_+$ と $D_-$ の平面座標はいずれも $x\in E\subset\mathbb P(B)$ とする。

\begin{theorem}[平滑化切断の全射性]
\label{thm:full-t1-v30}
三つの交代形式を任意に一次摂動する空間を
\[
 \mathcal V=(\Lambda^2(T\oplus B)^*)^{\oplus3}
\]
とする。局所平滑化切断を取る線形写像
\begin{equation}
 \mathcal V\longrightarrow
 H^0(E,H^3\otimes\theta)\oplus
 H^0(E,H^3\otimes\theta^{-1})
 \label{eq:full-t1-v30}
\end{equation}
は全射である。従って、$\Omega$ を固定したまま、一般の線形摂動の
十八個の節点の位置は
\[
 |H^3\otimes\theta|\times|H^3\otimes\theta^{-1}|
 \simeq\mathbb P^8\times\mathbb P^8
\]
を支配する。
\end{theorem}

\begin{proof}
$D_+$ の近くで使った座標
\[
 (u,Au,Cu),\qquad(0,d,\ell),\qquad
 \ell=(\ell_1,\ell_2,\ell_3)
\]
を用いる。$U_+$--$B$ 混合ブロックだけを摂動し、
第 $i$ 形式のそのブロックを
$\beta_i\in U_+^*\otimes B^*$ とする。
$U_+\times U_+$ 上ではこの摂動は零なので、
$A$ を消去する一次線形作用素は変わらない。
最後の混合方程式は、一次まで
\[
 M_i d+C^t\varepsilon_i\ell+t\beta_i(-,\ell)=0
\]
となる。$\sum_i\ell_i\varepsilon_i\ell=0$ を用いて $C$ を消去すると、
\begin{equation}
 M(\ell)d+tq(\ell)=0,\qquad
 q(\ell)=\sum_i\ell_i\beta_i(-,\ell)
 \label{eq:mixed-quadratic-v30}
\end{equation}
を得る。$\beta_i$ を動かすと $q$ は任意の
$U_+^*$ 値二次形式になる。これは各座標について
$\ell_i\ell_j$ の全てが現れることから直ちに従う。

平面 $\Pi_+$ に沿う $d$ の束は $\mathcal O(1)^3$ であり、
式\eqref{eq:mixed-quadratic-v30}の方程式の束は $\mathcal O(2)^3$ である。
局所平滑化切断は $q$ の余核での像に等しい。
行列式表示の分解を $\mathcal O(3)$ で捻ると
\begin{equation}
 0\longrightarrow\mathcal O_{\mathbb P^2}(1)^3
 \xrightarrow{M}\mathcal O_{\mathbb P^2}(2)^3
 \longrightarrow i_*(H^3\otimes\theta)\longrightarrow0
 \label{eq:t1-resolution-v30}
\end{equation}
となる。$H^1(\mathbb P^2,\mathcal O(1))=0$ なので
$H^0(\mathcal O(2)^3)\to H^0(H^3\otimes\theta)$ は全射である。

$D_-$ では転置行列を使い、$U_-$--$B$ 混合ブロックが
$H^0(H^3\otimes\theta^{-1})$ に全射する。
前者の混合ブロックは $U_-\oplus\ell$ 上で零であり、後者も
$U_+\oplus\ell$ 上で零であるから、二つの写像は独立である。
良い線形摂動は $\mathcal V$ の空でない開集合をなす。
全射線形写像は開写像なので、その像は二つの切断空間の稠密開集合を含む。
切断のスカラー倍を除けば最後の主張を得る。
\end{proof}

\begin{remark}[切断をあらかじめ指定した場合]
\label{rem:prescribed-zeros-v30}
任意に指定した二つの切断がそれぞれ単純零点のみを持つとき、
それらを一次項として持つ摂動を選び、さらに一般の二次項を加えれば、
指定した十八点を持つ平滑化を作れる。
実際、$s(t)=s_0+t s_1+t^2s_2$ と置くと、固定した $t\ne0$ において
$s_2$ から全切断空間への評価は全射である。
従って一般の $s_2$ について一般ファイバーは滑らかであり、
必要なら原点の近傍を縮める。
中央ファイバーの各零点の局所式は依然として $ab=tc+$ 高次項なので
通常二重点である。この主張は、任意に指定した切断が既に
線形の一変数摂動だけで実現するという主張ではない。
\end{remark}

\subsection{原点、平面座標、および線束の符号}

以後 $H\simeq\mathcal O_E(3o)$ を満たす点 $o\in E$ を一つ選び、
$o$ を原点とする。固定した $H$ に対してこの選択は九通りあり、
族の上では有限エタール被覆を取れば選べる。
同様に、単純な九点零因子をそれぞれ順序付けるため、
必要な有限エタール被覆を取る。
この順序付けが例外曲線 $E_i^\pm$ と整数格子のマーキングを定める。
原点を用いて
\[
 E\simeq\operatorname{Pic}^0(E),\qquad
 a\longleftrightarrow\theta=\mathcal O_E(a-o)
\]
と同一視し、以下の和は $E$ の群法則によるものとする。

元の平面 $\Pi_\pm$ 上の零点を $q_i^\pm\in E$ とする。
その線形同値条件は
\begin{equation}
 \sum_{i=1}^9q_i^+=a,\qquad
 \sum_{i=1}^9q_i^-=-a.
 \label{eq:sum-points-v30}
\end{equation}
線織面の底曲線の座標を $y$ とすると、既に得た二つの貼合せ写像は
$y=x+a$ と $y=x-a$ である。従って、共通の底曲線上の例外曲線の交点は
\begin{equation}
 P_i^+=q_i^++a,\qquad P_i^-=q_i^--a.
 \label{eq:translated-points-v30}
\end{equation}
両端の平面直線束をこの同じ $E$ 上に移したものを
$\mathscr H_+,\mathscr H_-$ と書く。
次数 $d$ の束 $A$ に対して $t_a^*A=A\otimes\theta^{-d}$ なので、
\begin{equation}
 \mathscr H_+=H\otimes\theta^3,\qquad
 \mathscr H_-=H\otimes\theta^{-3},
 \qquad
 \mathcal O_E\Bigl(\sum_iP_i^\pm\Bigr)
     =H^3\otimes\theta^{\pm10}.
 \label{eq:translated-bundles-v30}
\end{equation}
特に
$N_{D_+/Y_+}=\theta^{-1}$、$N_{D_-/Y_-}=\theta$ となり、
中間成分の法束 $\theta,\theta^{-1}$ と打ち消し合う。

$R$ 上で $S=D_+$ を選ぶと
\[
 \mathcal O_R(S)|_{D_+}=\theta,\qquad
 \mathcal O_R(S)|_{D_-}=\mathcal O_E.
\]
従ってPl\"ucker束の中間成分への制限は数値同値より強く
\begin{equation}
 \mathscr H_R\simeq
 \mathcal O_R(6S)\otimes p^*(H\otimes\theta^{-3})
 \label{eq:plucker-exact-v30}
\end{equation}
と決まる。この束の二つの制限はそれぞれ
$\mathscr H_+$ と $\mathscr H_-$ である。

\subsection{明示的な貼合せperiod}

前節の記号
\[
 \Lambda=I_{1,9}\oplus U\oplus I_{1,9},\qquad
 j_+=(d_+,-S,0),\quad j_-=(0,S,-d_-),\quad
 L=\langle j_+,j_-\rangle^\perp/\langle j_+,j_-\rangle
\]
を用いる。代表元を
\begin{equation}
 x=\left(AH_+-\sum_i b_iE_i^+,\ cS+dF,\
          A'H_--\sum_i b_i'E_i^-\right)
 \label{eq:period-coordinates-v30}
\end{equation}
と書く。$x\in J^\perp$ は
\begin{equation}
 \sum_i b_i=3A-d,\qquad \sum_i b_i'=3A'-d
 \label{eq:period-degree-constraints-v30}
\end{equation}
と同値である。

\begin{theorem}[periodの公式と貼合せ判定]
\label{thm:period-formula-v30}
上の順序と符号を固定すると、Type~IIの貼合せperiodは
\begin{align}
 \psi(x)
 &=\mathscr H_+^{A}\otimes\mathscr H_-^{-A'}\otimes\theta^{-c}\notag\\
 &\qquad\otimes\mathcal O_E\left(-\sum_i b_iP_i^+
                                  +\sum_i b_i'P_i^-\right)
                                                            \label{eq:period-multiplicative-v30}\\
 &=-\sum_i b_iq_i^+ +\sum_i b_i'q_i^- +(2d-c)a
      \quad\text{（$E$ 上の加法表示）}                      \label{eq:period-additive-v30}
\end{align}
で与えられる。これは $L\to\operatorname{Pic}^0(E)$ の準同型であり、
$\psi(h)=0$ を満たす。
また、数値類 $x$ を持つ成分上の線束を $\widetilde X$ 上の一つの線束に
貼り合わせられるための必要十分条件は $\psi(x)=0$ である。
\end{theorem}

\begin{proof}
両端の線束を $\mathscr L_+,\mathscr L_-$ とし、
中間成分の線束を
$\mathscr L_R=\mathcal O_R(cS)\otimes p^*B$、$\deg B=d$ と書く。
二つの貼合せの不一致の比
\begin{equation}
 \mathscr L_+|_E\otimes(\mathscr L_R|_{D_+})^{-1}
       \otimes\mathscr L_R|_{D_-}\otimes(\mathscr L_-|_E)^{-1}
 \label{eq:ratio-obstruction-v30}
\end{equation}
は $B$ に依存せず、式\eqref{eq:period-multiplicative-v30}に等しい。
次数条件\eqref{eq:period-degree-constraints-v30}からその次数は零である。
$B$ の次数零部分を動かして一方の不一致を零にでき、
他方も零になる条件がまさにこの比の自明性である。
三重点がないため、これで線束の貼合せに必要な条件は尽くされる。

$j_+$ と $j_-$ に対してこの比はそれぞれ
$N_{D_+/Y_+}\otimes\theta$ と
$N_{D_-/Y_-}\otimes\theta^{-1}$ なので自明であり、$L$ 上に降りる。
$h=(H_+,6S+3F,H_-)$ に対しては
$\mathscr H_+\mathscr H_-^{-1}\theta^{-6}=\mathcal O_E$ となる。
最後に式\eqref{eq:translated-points-v30},
\eqref{eq:translated-bundles-v30}を代入する。
$a$ の係数は
\[
 3(A+A')-c-\sum_i b_i-\sum_i b_i'=2d-c
\]
となり、加法表示を得る。
\end{proof}

\begin{remark}[混合Hodge構造との関係]
\label{rem:period-mhs-v30}
二つの有理成分の場合に、式\eqref{eq:ratio-obstruction-v30}の制限束の比を
periodとする定義は \cite[\S2.8]{KondoTypeII} にあり、
それがCarlsonの混合Hodge構造の拡大類と一致することは
同論文\S2.9(i)に明記されている。
三成分の場合も、$x$ に $c j_+$ を加えて中間成分の $S$ 係数を零にすると、
両端の組
\[
 \bigl(\mathscr L_+\otimes\mathcal O_{Y_+}(cD_+),\mathscr L_-\bigr)
\]
に対する同じ比になる。これは線織成分をファイバー方向に縮約する際の
整数格子の同一視とも整合する。
従って本節で計算するのは中央ファイバーの重み二以下の拡大類、
およびその極限混合Hodge構造への対応部分である。
平滑化パラメータの正規化や、重み三からの追加の拡大データまで
この式で決まるとは主張しない。
\end{remark}

\begin{corollary}[根のperiodとねじれ方向]
\label{cor:root-period-v30}
偏極直交格子 $L_h=h^\perp$ の根について
\begin{equation}
 \psi(E_i^+-E_j^+)=q_i^+-q_j^+,\qquad
 \psi(E_i^--E_j^-)=q_j^--q_i^-
 \label{eq:root-period-v30}
\end{equation}
が成立する。従って各端の九点が相異なる条件は、
$L_h$ の全ての根のperiodが零でない条件と同値である。
また
\[
 f=(d_+,0,0),\qquad
 v=(3H_+,-20S+9F,3H_-)
\]
に対し
\begin{equation}
 \psi(f)=-a,\qquad \psi(v)=38a
 \quad\text{すなわち}\quad
 \psi(v)=\theta^{38}
 \label{eq:theta38-v30}
\end{equation}
となる。$\theta$ の位数が $m<\infty$ ならば、$nv$ が貼合せ可能になる
最小の正整数は $m/\gcd(m,38)$ である。
$\theta$ が無限位数なら、非零の整数倍 $nv$ は貼合せ可能でない。
\end{corollary}

\begin{proof}
公式への代入による。根が二つの $A_8$ 系の根のみであることは
前節で整数格子として示した。$v$ では $b_i=b_i'=0$、$d=9$、$c=-20$
なので $2d-c=38$ である。
なお、整数格子内でさらに
\[
 v=3h-38f\quad\text{in }L
 \qquad
 (v-3h+38f=38j_+\quad\text{in }\Lambda)
\]
が成立する。従って指数三十八はこの整数関係からも直接読み取れる。
\end{proof}

\subsection{偏極付きperiodへの次数九のisogeny}

固定した $(E,o)$ に対して
\begin{equation}
 \mathcal P_E=
 \left\{(q_1^+,\ldots,q_9^+,q_1^-,\ldots,q_9^-)\in E^{18}:
         \sum_iq_i^++\sum_iq_i^-=0\right\}
 \simeq E^{17}
 \label{eq:point-abelian-v30}
\end{equation}
と置く。$a=\sum_iq_i^+$ を回復すれば、式\eqref{eq:sum-points-v30}は
自動的に成り立つ。条件 $2a\ne0$ と各九点の相異なりは
$\mathcal P_E$ のZariski開条件である。
また、$a$ を固定したファイバーは次元十六であり、
順序を忘れれば前述の二つの完全線形系に等しい。

\begin{theorem}[明示的な偏極periodの被覆]
\label{thm:period-isogeny-v30}
式\eqref{eq:period-additive-v30}はisogeny
\begin{equation}
 \Phi_E:\mathcal P_E\longrightarrow
 \operatorname{Hom}(L/\mathbb Zh,E)
 \label{eq:period-isogeny-v30}
\end{equation}
を定める。その核は
\[
 \Delta E[3]=\{(b,\ldots,b):b\in E[3]\}
\]
であり、次数は九である。特に $\Phi_E$ はエタールで、
\[
 \mathcal P_E/\Delta E[3]\simeq\operatorname{Hom}(L/\mathbb Zh,E).
\]
\end{theorem}

\begin{proof}
自由アーベル群
\[
 P=\mathbb Z^{18}/\mathbb Z(1,\ldots,1)
\]
を用いると $\mathcal P_E=\operatorname{Hom}(P,E)$ である。
periodの整数係数写像は
\begin{equation}
 \sigma:J^\perp\longrightarrow P,\qquad
 x\longmapsto
 (2d-c-b_1,\ldots,2d-c-b_9,b_1',\ldots,b_9')
       \pmod{\mathbb Z(1,\ldots,1)}.
 \label{eq:period-character-v30}
\end{equation}
既に $J$ と $h$ は核に含まれる。

まず核を求める。係数ベクトルが全て同じ整数 $k$ のとき
$b_i=2d-c-k$、$b_i'=k$ であり、次数条件から
\[
 3A=19d-9c-9k,\qquad 3A'=d+9k
\]
となる。従って $d=3r$ と書ける。
直接計算すると
\[
 x-rh=-(c-6r+k)j_+-k j_-.
\]
従って、$L$ 上での核は正確に $\mathbb Zh$ である。

次に像を求める。係数の総和は
\[
 18d-9c-3A+3A'
\]
であり、三の倍数である。対角ベクトルを足すと総和は十八の倍数だけ
変わるので、この合同条件は $P$ 上で意味を持つ。
逆に $(u_1,\ldots,u_9,w_1,\ldots,w_9)$ が総和三の倍数なら、
\[
 d\equiv-\sum_iw_i\equiv\sum_i u_i\pmod3
\]
を満たす整数 $d$ を選び、任意の整数 $c$ に対し
\[
 b_i=2d-c-u_i,\quad b_i'=w_i,\quad
 A=\frac{19d-9c-\sum_i u_i}{3},\quad
 A'=\frac{d+\sum_iw_i}{3}
\]
とすれば、そのベクトルは $\sigma(x)$ になる。
従って
\begin{equation}
 0\longrightarrow L/\mathbb Zh\xrightarrow{\sigma}P
 \xrightarrow{\text{係数総和 mod }3}\mathbb Z/3\longrightarrow0
 \label{eq:period-index-three-v30}
\end{equation}
は完全である。

これに $\operatorname{Hom}(-,E)$ を適用する。
$[3]:E\to E$ が全射なので、得られる写像は全射であり、その核は
$\operatorname{Hom}(\mathbb Z/3,E)=E[3]$ である。
総和写像に対応する埋込みは全ての座標への同時平行移動なので、
核は $\Delta E[3]$ である。標数零であるからこの核はエタールで九点を持つ。
\end{proof}

\begin{corollary}[良い配置のperiod開集合]
\label{cor:period-open-v30}
$\mathcal P_E^\circ$ を $2a\ne0$ かつ各端の九点が相異なる配置の開集合とし、
\[
 \mathcal B_E^\circ=
 \{\psi\in\operatorname{Hom}(L/\mathbb Zh,E):
      2\psi(f)\ne0,\quad
      \psi(\alpha)\ne0\ \text{for every root }\alpha\in L_h\}
\]
と置く。このとき
\[
 \mathcal P_E^\circ/\Delta E[3]\simeq\mathcal B_E^\circ.
\]
従って根のperiodが零となる部分は、配置モデルでは各端の点の衝突に正確に対応する。
\end{corollary}

\begin{proof}
系\ref{cor:root-period-v30}により
$\mathcal P_E^\circ=\Phi_E^{-1}(\mathcal B_E^\circ)$ である。
定理\ref{thm:period-isogeny-v30}をこの開集合に制限すればよい。
\end{proof}

\begin{remark}[原点変更と残る有限対称性]
$H=\mathcal O_E(3o)$ を満たす原点を $o+b$、$b\in E[3]$ に変更すると、
同じ幾何学的点の群座標は一斉に $-b$ だけ変わる。
一方、原点を固定して全ての点を $b\in E[3]$ だけ動かす操作は、
$H$ を保つ $E$ の平行移動であり、平面三次曲線の射影自己同型に延長される。
両端にこの自己同型を同時に作用させると、$a$ は
$9b=0$ により変わらず、貼合せも保たれる。
従って次数九の核には、配置座標に残る具体的な同時平行移動としての意味がある。
九点の置換、$+$ と $-$ の交換、特別な $E$ の自己同型は別の有限作用であり、
本定理では商を取らず、マーキングにより固定している。
\end{remark}

\begin{remark}[直交格子への制限]
\label{rem:period-index38-v30}
$L$ はユニモジュラーで $h$ は原始的、$h^2=38$ なので
\[
 0\longrightarrow L_h\longrightarrow L/\mathbb Zh
       \longrightarrow\mathbb Z/38\longrightarrow0
\]
は完全である。従って偏極periodを $L_h$ に制限するisogenyの核は
$E[38]$、次数は $38^2$ である。
式\eqref{eq:period-isogeny-v30}との合成の整数係数写像の指数は
$3\cdot38=114$、isogenyの次数は $114^2=12996$ となる。
根格子 $A=A_8(-1)^2$ と $v$ に対して既に得た
$[L_h:A\oplus\mathbb Zv]=27$ とも整合する。実際、$P$ における
$\sigma(A\oplus\mathbb Zv)$ の指数は $3078=27\cdot114$ である。
この有限指数の違いを無視して $\operatorname{Hom}(L_h,E)$ を使うと、
偏極を零にするperiodの整数構造を取り違える。
\end{remark}

\subsection{逆Gram族の境界period支配性とHalphen指数}

\begin{theorem}[元の族の境界period支配性]
\label{thm:boundary-period-dominance-v30}
逆Gramパラメータと一般の平滑化方向を共に動かし、
両端と九点を順序付け、必要な局所マーキングを取る。
そのType~II period写像は、固定した整数境界型 $(L,h)$ に付随する
十八次元の偏極period族
\[
 \{(E,\psi):\psi\in\operatorname{Hom}(L/\mathbb Zh,E)\}
\]
を支配する。従ってその有限・算術群による標準的な商においても、
同じ境界成分の一般periodを与える。
これは個々の逆Gramパラメータを固定したまま十八次元を動かすという
主張ではない。固定した $(E,\theta)$ では九点配置が与える十六方向、
$\theta$ の変化が与える一方向、$E$ の変化が与える一方向を合わせた結論である。
\end{theorem}

\begin{proof}
前節までに証明した逆Gram族の支配性から、$\Omega$ を動かした
$(E,\{\theta,\theta^{-1}\})$ の像は、これらの組の二次元モジュライの
稠密開集合を含む。両端を順序付ける二重被覆を取り、
必要なら $\theta$ を $\theta^{-1}$ に取り替えると、
$(E,\theta)$ の稠密開集合を得る。
固定した一般の $\Omega$ に対して、定理\ref{thm:full-t1-v30}により
二つの完全線形系の積を支配する。
点を順序付けた被覆を合わせると、得られる $(E,q_i^\pm)$ は
式\eqref{eq:point-abelian-v30}の族の稠密開集合を支配する。
定理\ref{thm:period-isogeny-v30}は各 $E$ 上でエタール全射であり、
原点の選択後に族としても成立するので、period族への支配性が従う。
\end{proof}

\begin{remark}[モジュライに関して証明した範囲]
本定理は既知のType~II period空間を利用する支配性の主張である。
中央ファイバーの同型類のfine moduli空間を新たに構成したという主張でも、
別の境界型を含む全ての境界成分を覆うという主張でもない。
また、根のperiodが零となる超曲面は配置空間では同じ端の点の衝突対角線であるが、
この事実だけから任意に選んだ平滑化の一般ファイバーに対応する有効根や
Picard数の上昇が生じるとは結論しない。
\end{remark}

\begin{corollary}[十分大きい全てのHalphen指数]
\label{cor:cofinite-halphen-v30}
ある滑らかな平面三次曲線 $E_0$ と整数 $m_0$ が存在し、
全ての $m\ge m_0$ に対して、元の逆Gram族の良いパラメータと平滑化を選び、
底曲線を $E_0$ としたまま二つの端成分をともに指数 $m$ のHalphen曲面にできる。
\end{corollary}

\begin{proof}
前の証明における $(E,\theta)$ の稠密開像から、
ある固定した $E_0$ 上で、有限個の $\theta$ を除く全てが良い
逆Gramパラメータにより実現するように $E_0$ を選べる。
除外された有限集合に属する有限位数点の位数には上限がある。
それを超える全ての $m$ について、位数が正確に $m$ の点を選ぶことができる。
二つの零因子を一般に選ぶと、端成分の反標準法束は
$\theta^{-1}$ と $\theta$ である。
前節の反標準束の判定から、両端のHalphen指数は正確に $m$ となる。
\end{proof}

この系は存在的な結果であり、$m_0$ の有効上界は与えない。
元の制約付き逆Gram族で全ての $m\ge3$ が実現するかどうかは、
引き続き別の問題である。一方、period空間上ではこの問題の方向は明瞭であり、
Halphen指数 $m$ の条件は $\psi(f)$ が正確に位数 $m$ を持つ条件として書ける。
また $v$ の貼合せ条件は $38\psi(f)=0$ である。

\begin{corollary}[貼合せ類の跳躍とHalphen層]
\label{cor:cartier-strata-v30}
固定した $E$ のperiod空間において、$f\bmod h$ の評価写像
\[
 \operatorname{Hom}(L/\mathbb Zh,E)\longrightarrow E,
 \qquad\psi\longmapsto\psi(f)
\]
は連結な十六次元核を持つ。従って正確なHalphen指数 $m\ge3$ の層は、
良い配置の開条件を課す前には
\[
 J_2(m)=m^2\prod_{p\mid m}(1-p^{-2})
\]
個の連結成分からなる余次元一の層である。
また、複素数体上の非常に一般の中央ファイバーにおいて
$\ker\psi=\mathbb Zh$ であり、そのPicard群の階数は三である。
\end{corollary}

\begin{proof}
前節の双曲平面の基底 $e,f$ と $E_8$ 基底の元
$\alpha_8^+=H_+-E_1^+-E_2^+-E_3^+$ に対し
\[
 e-9\alpha_8^+\in L_h,\qquad f\cdot(e-9\alpha_8^+)=1
\]
である。従って $f\bmod h$ は $L/\mathbb Zh$ 内で原始的であり、
評価写像は基底を選べば $E^{17}\to E$ の一つの射影となる。
正確に位数 $m$ の点の個数は $J_2(m)$ なので前半を得る。

任意の非零 $\lambda\in L/\mathbb Zh$ に対して
$\{\psi:\psi(\lambda)=0\}$ は真の閉部分集合である。
従って可算個のこれらの部分集合を除けば $\ker\psi=\mathbb Zh$ となる。
period支配性により、元の族でも非常に一般の点にこの条件が成立する。
線束の貼合せ判定から、中央ファイバーの数値Cartier類は
$J^\perp$ における $\mathbb Zh$ の逆像であり、階数は $2+1=3$ である。
実際のPicard群からこの数値群への写像も単射である。
数値的に零の束は両端で自明であり、中間成分では $p^*B$、$B\in\operatorname{Pic}^0(E)$
となるが、二重曲線での貼合せから $B$ も自明になる。
二重曲線の双対グラフは木なので、成分上で自明な線束の貼合せにも
非自明な自由度は残らない。
\end{proof}

ここでの階数三は中央ファイバーのPicard群についての主張であり、
平滑な一般ファイバーのPicard数を三とする主張ではない。
二つの垂直方向は $L$ への商で消える。

\subsection{再現可能な整数計算と新規性の範囲}

添付の \texttt{period\_lattice\_v30.py} は、二十二次元の成分格子から
$L$ のユニモジュラー基底を作り直し、式\eqref{eq:period-character-v30}の
Smith標準形を整数演算で求める。
偏極periodの係数写像は非零対角成分 $(1^{16},3)$、
$L_h$ に制限した写像は $(1^{16},114)$ となる。
さらに $\psi(h)=0$、$\psi(v)=38a$、根格子を含む指数二十七との整合、
混合摂動から三つの二次形式への全射性を検査する。
証明自体は上の議論で完結しており、この計算は整数の符号と指数の独立な確認である。

一般のType~II period、Carlson拡大、periodのTorelli理論を本研究の
新しい結果として数えない。本節で得た具体的な内容は、
制約付き逆Gram族のGrassmann平滑化が二つの完全線形系を独立に動かすこと、
その幾何学的配置が次数九の明示的isogenyによって偏極periodを与えること、
およびその結果として元の族がこの境界型の全period方向を支配することである。

% 挿入位置: v25 の補助Gram曲線との二重交差命題の後。
% この節の計算は同一Hensel分枝を使い、点配置のmonodromyを主張しない。
\section{補助Gram曲線の通常二重点と既約性}
\label{sec:complementary-nodes-v26}
\label{subsec:complementary-gram-nodes-v26}

補助Gram曲線の固定点特異点は、全incidence方程式を用いると
通常二重点まで決定できる。この局所的結論には、新しい一般性仮定や計算を要しない。

\begin{proposition}[固定点における非退化な法線形式]
\label{prop:fixed-gram-node-v26}
第\ref{sec:prym-setup-v26}節の全incidence、可逆対称pairing、
平面全体での $\operatorname{rank}P=p$、有限被約な $Z$ の条件を満たすとする。
各固定点 $x_0$ において、$L(x_0)$ の行空間上の $A$ による対称形式は非退化である。
従って $P(x_0)AP(x_0)^t$ の階数は $p$、
$B(x_0)AB(x_0)^t$ の階数は $q-1$ であり、
$x_0$ は $\Gamma_B$ の通常二重点である。
\end{proposition}

\begin{proof}
平面の座標を変え、$x_0=[1:0:0]$ とする。
以下では $M=L_0=L(x_0)$、$L(1,s,t)=M+sL_1+tL_2$ と書く。
被約な孤立退化点では $\operatorname{rank}M=5$ である。
実際、階数が四以下なら六次小行列式の全ての一次微分が零となり、
局所的な退化イデアルは極大イデアルを生成できない。
左核の生成元を行ベクトル $u=(0,b)$ と取る。
固定点なので $uJ=-u$ であり、$uM=0$ である。

$w_i=uL_i=bB_i$ $(i=1,2)$ と置く。
全incidence方程式を左から $u$ で掛けると
\[
 uL_iAM^tJ=uJMAL_i^t=0
\]
となるので、
\begin{equation}\label{eq:fixed-normal-orthogonality-v26}
 w_iAM^t=0,\qquad Aw_i^t\in K:=\ker M
\end{equation}
を得る。ここで $\dim K=2$ である。

階数五の行列の退化locusに対する法線微分は
\[
 T_{x_0}\mathbb P^2\longrightarrow K^*,\qquad
 (a_1,a_2)\longmapsto (a_1w_1+a_2w_2)|_K
\]
である。局所的に五次可逆小行列を消去すれば、退化イデアルは残る一行二列の
二成分で生成されるので、$Z$ が被約であることから、この微分は同型となる。
従って $w_1,w_2$ は独立であり、式\eqref{eq:fixed-normal-orthogonality-v26}により
$Aw_1^t,Aw_2^t$ は $K$ の基底をなす。
この基底に関する法線微分の行列は
\begin{equation}\label{eq:fixed-normal-gram-v26}
 H_0=(w_iAw_j^t)_{1\leq i,j\leq2}
\end{equation}
なので、$H_0$ は非退化である。

行空間 $R=\operatorname{Row}M$ の直交補空間は
$\langle w_1,w_2\rangle$ であり、そこに制限した形式が
式\eqref{eq:fixed-normal-gram-v26}である。
従って $R$ も非退化である。
さらに
\[
 R=\operatorname{Row}P(x_0)\mathbin{\perp}\operatorname{Row}B(x_0)
\]
は、次元 $p$ と $q-1$ の直交直和であるから、両項の制限形式も非退化となる。

$B$ の行の定数基底変換により、最後の行を $bB$ とし、
他の $q-1$ 行を $C(s,t)$ とする。
$C(0,0)AC(0,0)^t$ は可逆であり、最後の行は $sw_1+tw_2$ である。
式\eqref{eq:fixed-normal-orthogonality-v26}により、
Gram行列の混合ブロックは $(s,t)$ の二次以上の項からなる。
Schur補行列を取ると、$\det(BAB^t)$ の二次初項は
\[
 \det\bigl(C(0,0)AC(0,0)^t\bigr)
 (s,t)H_0(s,t)^t
\]
となる。これは非退化な二次形式であるから、$x_0$ は通常二重点である。
\end{proof}

\begin{lemma}[一つの次数による射影的な長さの上界]
\label{lem:jacobian-length-bound-v26}
$k$ を無限体、$S=k[x,y,z]$、$I\subset S$ を斉次イデアルとする。
ある $D$ と $f\leq D$ に対し $\dim_k(S/I)_D\leq f$ なら、
$\operatorname{Proj}(S/I)$ は零次元または空であり、その長さは $f$ 以下である。
\end{lemma}

\begin{proof}
$I$ を飽和しても指定次数の商空間の次元は増えない。
$R=S/I^{\mathrm{sat}}$ とし、$R$ が零でなければ、
飽和性と $k$ の無限性により $R$ 上の非零因子となる一次式 $\ell$ を取れる。
$h_d=\dim_k R_d$、$h_{-1}=0$ と置くと
\[
 h_d-h_{d-1}=\dim_k(R/\ell R)_d\geq0.
\]
$R/\ell R$ は次数一で生成されるので、そのある次数の成分が零なら、
以後の成分も全て零である。
もし次数零から $D$ までが全て非零なら $h_D\geq D+1$ となり、
$h_D\leq f\leq D$ に反する。
従って $h_d$ は $D$ 以前に定常となり、その定常値は $f$ 以下である。
これは射影スキームが零次元または空で、その長さが $f$ 以下であることを意味する。
\end{proof}

\begin{theorem}[補助曲線の特異点の完全な決定]
\label{thm:complementary-gram-nodal-v26}
固定成分 $Z_1,Z_3$ は、定理\ref{thm:gram-smooth-v25}の条件に加えて
次の性質を持つ非空開部分を含む。
補助Gram曲線 $\Gamma_B$ は幾何学的に既約であり、
その特異点は指定反対合の固定点ちょうど全部で、全て通常二重点である。
\[
\begin{array}{c|c|c|c|c}
 \text{成分}&\deg C&\deg\Gamma_B&
 \Gamma_B\text{ の通常二重点数}&\text{正規化の種数}\\\hline
 Z_1&6&6&3&7\\
 Z_3&4&8&5&16
\end{array}
\]
非固定点では $C$ と $\Gamma_B$ はともに滑らかであり、
各点の局所交差数は二である。
従って $Z_1$ では十八点、$Z_3$ では十六点で通常の二重接触を持つ。
\end{theorem}

\begin{proof}
定理\ref{thm:gram-smooth-v25}と同じ代数的 $101$ 進Hensel点を使う。
基点、pivot変数、自由変数、seedは、被約二十一点配置の既存証明書と全て照合する。
$G=\det(BAB^t)$、$d=2q$、$D=3d-5$ とし、
$G_x,G_y,G_z$ に次数 $D-d+1$ の全単項式を掛け、
次数 $D$ の単項式基底で並べたMacaulay行列を作る。
指定した行と列による正方小行列について、次の値を得る。
\[
\begin{array}{c|c|c|c|c|c|c}
 \text{成分}&d&D&\text{行列の大きさ}&
 \text{小行列の次数}&v_{101}(\det)&\text{保証精度}\\\hline
 Z_1&6&13&135\times105&102&48&256\\
 Z_3&8&19&273\times210&205&63&256
\end{array}
\]
行列の全成分と小行列式はHensel座標の整係数多項式であり、分母を導入しない。
従って入力精度 $101^{256}$ が小行列式にもそのまま保証される。
表の付値は保証精度未満なので、真の標数零のHensel点で小行列式は非零である。

$J_G=(G_x,G_y,G_z)$ とすると、表から
\[
 \dim(S/J_G)_{13}\leq3\quad(Z_1),\qquad
 \dim(S/J_G)_{19}\leq5\quad(Z_3)
\]
を得る。標数零なのでEuler公式により $G\in J_G$ であり、
$\operatorname{Proj}(S/J_G)$ は $\Gamma_B$ のJacobian特異スキームである。
補題\ref{lem:jacobian-length-bound-v26}により、その長さはそれぞれ三、五以下である。
一方、命題\ref{prop:fixed-gram-node-v26}により、既知の相異なる三、五個の固定点は
通常二重点であり、その各点におけるJacobian特異スキームの長さは一である。
従って特異スキームはこれらの被約固定点だけに一致する。

特異点が有限なので $\Gamma_B$ は幾何学的に被約である。
もし次数 $a$ と $d-a$ の二つの正次数曲線に分けられれば、
B\'ezoutの定理により、その交差数は
\[
 a(d-a)\geq d-1
\]
である。全特異点が通常二重点なので、異なる成分の交点は横断的で、
各点の交差数は一である。
従って少なくとも $d-1$ 個の特異点が必要になるが、
ここでは $3<5$、$5<7$ なので矛盾する。
よって $\Gamma_B$ は幾何学的に既約である。
正規化の種数は平面曲線の算術種数から通常二重点数を引き、
$10-3=7$、$21-5=16$ となる。

以上の非消滅条件は、既知の固定点が相異なる通常二重点である開部分で開条件を定める。
同じHensel点が小さいGram曲線の滑らかさも満たすので、
元の固定成分内に求める非空開部分を得る。
最後に、固定点は $C$ の外にあり、これ以外に $\Gamma_B$ の特異点はない。
命題\ref{prop:complementary-gram-contact-v25}の
$\Gamma_B|_C=2D_{\rm mov}$ により、
非固定点での局所交差数は二となる。
\end{proof}

\paragraph{再現入口.}
\nolinkurl{complementary_gram_nodes_v26.py} と
\nolinkurl{complementary_gram_nodes_v26_certificate.json} が、
小行列の行番号・列番号、行列式の付値と正規化した剰余、Gram多項式のhash、
既存Hensel証明書との同一性を記録する。
既定の実行は読取り専用照合である。
この計算は標数零の特異スキームの上界を証明するものであり、
有限体の算術Galois群から幾何学的monodromyを推定するものではない。

\begin{remark}[補助曲線の対称行列表示も線束を定める]
\label{rem:complementary-nodal-prym-v26}
上の開部分では $Q_B=BAB^t$ の階数は $\Gamma_B$ の全点で $q-1$ である。
固定点では命題\ref{prop:fixed-gram-node-v26}による。
それ以外で階数が $q-2$ 以下なら随伴行列が零となり、余分な特異点が生じてしまう。
従って
\[
 0\longrightarrow\mathcal O_{\mathbb P^2}(-2)^q
 \xrightarrow{Q_B}\mathcal O_{\mathbb P^2}^q
 \longrightarrow j_*\mathcal L_B\longrightarrow0
\]
の余核は、節点を含めて $\Gamma_B$ 上の線束である。
Cartier因子に対する双対性と対称性を用いる同じ証明により、
\[
 \eta_B:=\mathcal L_B(1-q),\qquad
 \eta_B^{\otimes2}\simeq\mathcal O_{\Gamma_B},\qquad
 H^0(\Gamma_B,\eta_B)=0
\]
を得る。従って補助節点曲線にも非自明な二捩れ線束が付随し、
連結な有限エタール二重被覆を定める。
ここでは、この線束の正規化への引き戻しが非自明であることや、
小さいGram曲線の二重被覆との同一視までは主張しない。
\end{remark}


\section{検証と適用範囲}
同じ既存成分上の非空開条件、Gram写像の微分階数、Hermite--Pfaffian恒等式、
境界チャートの全式、整数格子の計算は、付属の厳密な有限証明書で再現できる。
線束・商写像・Kulikov解消、平滑化切断の全射性、周期の貼合せ判定と支配性は、本文の幾何学的証明による。
主な入口は \nolinkurl{gram_curves_v25.py}、\nolinkurl{gram_dominance_v26.py}、
\nolinkurl{hermite_pfaffian_v27.py}、\nolinkurl{grassmann_boundary_v28.py}、
\nolinkurl{ruled_moduli_dominance_v29.py}、\nolinkurl{kulikov_lattice_v29.py}、
\nolinkurl{period_lattice_v30.py} である。

ここでのKulikovモデルは一般の固定Gram表示と一般一次摂動に対する構成である。
同じ整数境界型の周期を支配するが、粗モジュライ上の相対コンパクト化や、その算術的降下を同時には主張しない。
計算する混合Hodge拡大は中央ファイバーの重み二以下の部分であり、全ての極限拡大データを決定したとはしない。
非捩れ性は幾何学的一般点および非常に一般の複素点についての結論である。
また $A_8\oplus A_8$ は偏極付き極限格子の根系であり、
一般K3ファイバーのPicard階数や、端成分上の有効 $(-2)$ 曲線の個数を主張するものではない。

\clearpage
\begingroup\small
\begin{thebibliography}{99}
\setlength{\itemsep}{2pt}
\setlength{\parskip}{0pt}
\bibitem{PaperII}
\emph{Fano72型関係代数の二十一点Torelliと反対代数}, 本研究の先行稿、第26版、2026年。
\nolinkurl{fano72_point_configurations_ja_v26.pdf}.

\bibitem{BeauvilleDeterminantal}
A.~Beauville, \emph{Determinantal hypersurfaces}, Michigan Math. J.\ \textbf{48} (2000), 39--64.
\url{https://arxiv.org/abs/math/9910030}.

\bibitem{BoehningBothmerPrym}
C.~B\"ohning and H.-C.~Graf von Bothmer,
\emph{On stable rationality of some conic bundles and moduli spaces of Prym curves},
Comment. Math. Helv.\ \textbf{93} (2018), 133--155.
\url{https://doi.org/10.4171/CMH/431}.

\bibitem{BazanMezzetti}
D.~Bazan and E.~Mezzetti, \emph{On the construction of some Buchsbaum varieties and the Hilbert scheme of elliptic scrolls in $\mathbb P^5$},
Geom. Dedicata \textbf{86} (2001), 191--204.
\url{https://doi.org/10.1023/A:1011926128564}.

\bibitem{QuarticAutomorphisms}
R.~Lercier, C.~Ritzenthaler, F.~Rovetta and J.~Sijsling,
\emph{Parametrizing the moduli space of curves and applications to smooth plane quartics over finite fields},
LMS J. Comput. Math.\ \textbf{17}, Special issue A (2014), 128--147.
\url{https://doi.org/10.1112/S146115701400031X}.

\bibitem{MukaiGenus20}
S.~Mukai, \emph{Polarized K3 surfaces of genus 18 and 20},
in Complex Projective Geometry, London Math. Soc. Lecture Note Ser.\ \textbf{179},
Cambridge University Press, 1992, 264--276.
\url{https://www.kurims.kyoto-u.ac.jp/~mukai/paper/Bergen.pdf}.

\bibitem{DolgachevCAG}
I.~V.~Dolgachev, \emph{Classical Algebraic Geometry: A Modern View},
Cambridge University Press, 2012.
著者公開原稿、\S3.2.2 and Exercise~3.12.
\url{https://sites.lsa.umich.edu/idolga/wp-content/uploads/sites/1334/2024/08/CAG.21.pdf}.
\url{https://doi.org/10.1017/CBO9781139084437}.

\bibitem{BruinPrymGenusThree}
N.~Bruin, \emph{The arithmetic of Prym varieties in genus 3},
Compos. Math.\ \textbf{144} (2008), 317--338.
\url{https://doi.org/10.1112/S0010437X07003314}.
\url{https://arxiv.org/abs/math/0408069}.

\bibitem{DolgachevRationality2008}
I.~V. Dolgachev,
\emph{Rationality of $\mathcal R_2$ and $\mathcal R_3$},
Pure Appl. Math. Q. \textbf{4} (2008), no.~2, 501--508.
\href{https://doi.org/10.4310/PAMQ.2008.v4.n2.a8}{doi:10.4310/PAMQ.2008.v4.n2.a8}.

\bibitem{PaperVI}
\emph{種数三Prymモジュライの相対有理性と二次被覆}, 本研究の公表稿VI、第30版、2026年。
\nolinkurl{prym_net_arithmetic_ja_v30.pdf}.
\bibitem{KulikovDegenerations}
V.~S.~Kulikov, \emph{Degenerations of K3 surfaces and Enriques surfaces},
Math. USSR-Izv. \textbf{11} (1977), no.~5, 957--989.
\url{https://doi.org/10.1070/IM1977v011n05ABEH001753}.
\bibitem{PerssonPinkhamDegenerations}
U.~Persson and H.~Pinkham, \emph{Degeneration of surfaces with trivial canonical bundle},
Ann. of Math. (2) \textbf{113} (1981), no.~1, 45--66.
\url{https://doi.org/10.2307/1971133}.


\bibitem{PaperVII}
\emph{六次Prym曲線の標準七次被覆と分岐曲線}, 本研究の公表稿VII、第30版、2026年。
\nolinkurl{sextic_prym_septic_covers_ja_v30.pdf}.
\bibitem{KondoTypeII}
S.~Kond\=o, \emph{Type II degenerations of K3 surfaces},
Nagoya Math. J. \textbf{99} (1985), 11--30.
\url{https://doi.org/10.1017/S0027763000021462}.
\bibitem{FriedmanTorelli}
R.~Friedman, \emph{A new proof of the global Torelli theorem for K3 surfaces},
Ann. of Math. (2) \textbf{120} (1984), no.~2, 237--269.
\url{https://doi.org/10.2307/2006942}.
\bibitem{AlexeevEngelHanAutomorphisms}
V.~Alexeev, P.~Engel and C.~Han,
\emph{Compact moduli of K3 surfaces with a nonsymplectic automorphism},
Trans. Amer. Math. Soc. Ser. B \textbf{11} (2024), 144--163.
\url{https://doi.org/10.1090/btran/136}.
本文の定理番号は著者公開稿（2022年9月13日）による。
\url{https://hanchangho.github.io/Papers/nonsympK3.pdf}.

\end{thebibliography}
\endgroup
\end{document}
